\documentclass{article}

\usepackage[preprint]{neurips_2026}

\usepackage[utf8]{inputenc}
\usepackage[T1]{fontenc}
\usepackage{hyperref}
\usepackage{url}
\usepackage{booktabs}
\usepackage{graphicx}
\usepackage{amsfonts}
\usepackage{array}
\usepackage{tabularx}
\usepackage{float}
\usepackage{nicefrac}
\usepackage{microtype}
\usepackage{xcolor}
\usepackage{enumitem}
\usepackage{fancyvrb}
\usepackage{listings}
\lstdefinestyle{kisspy}{
  language=Python,
  basicstyle=\ttfamily\scriptsize,
  keywordstyle=\color[rgb]{0.0,0.0,0.8}\bfseries,
  commentstyle=\color[rgb]{0.25,0.5,0.35}\itshape,
  stringstyle=\color[rgb]{0.73,0.13,0.13},
  showstringspaces=false,
  frame=single,
  framesep=2mm,
  numbers=left,
  numberstyle=\tiny\color{gray},
  numbersep=5pt,
  breaklines=true,
  tabsize=4,
  captionpos=b,
}
\lstnewenvironment{prompt}{%
  \lstset{language={},basicstyle=\ttfamily\scriptsize,frame=single,framesep=2mm,breaklines=true,columns=fullflexible}%
}{}

\title{KISS Sorcar: A Stupidly-Simple General-Purpose and Software Engineering AI Assistant}

\author{%
Koushik Sen \\ EECS Department, UC Berkeley\\ \texttt{ksen@berkeley.edu} \\ }

\begin{document}

\maketitle

\begin{center}
\textit{"Everything should be made as simple as possible, but not simpler." --- Albert Einstein}
\end{center}


\begin{abstract}
Large language models can generate code and call tools fluently, yet deploying them as practical software engineering assistants still exposes persistent gaps: finite context windows, a single mistake that can derail entire sessions, agents that get stuck in dead ends, AI slop, and generated changes that are difficult to review or revert.

We present \textbf{KISS Sorcar}, a general-purpose assistant and integrated development environment (IDE) built on top of the \textbf{KISS Agent Framework}, a stupidly-simple AI agent framework of roughly 2,500~lines of code for the core agents.  The framework addresses these gaps through a structured system prompt and a five-layer agent hierarchy in which each layer adds exactly one concern: budget-tracked ReAct execution, automatic continuation across sub-sessions via summarization, coding, and browser tools with parallel sub-agents, persistent multi-turn chat with history recall, and git worktree isolation so every task runs on its own branch.  Engineering principles are encoded in the agent's system prompt.

We implemented KISS Sorcar as a free, open-source, local-first, bring-your-own-key assistant that ships as a Visual Studio Code extension, a Claude-Code-style CLI (\texttt{sorcar}), and a browser/mobile web app.  Prompts and code are sent directly to the model provider or local endpoint the user configures, not through any intermediary servers.  The framework supports mixing models from multiple vendors in the same task (OpenAI, Anthropic, Gemini, Together, MiniMax, OpenRouter, Claude Code CLI, and Codex CLI), bundles a catalog of 501 models across 8 provider categories, includes 23 third-party messaging agents (Slack, Gmail, WhatsApp, SMS, phone-call control, and others), discovers Model-Context-Protocol (MCP) servers, loads Agent Skills, and supports browser automation, multimodal input, and Docker containers.

In this research, we deliberately prioritize output quality over speed:  giving a frontier model adequate time to validate its own output---running linters, type checkers, and tests---reduces the low-quality code common in faster but less thorough agents.  The entire system was built using itself in 4 months, providing a continuous stress test in which any bug was patched as it appeared.  On Terminal~Bench~2.0, KISS Sorcar achieves a 62.2\% overall pass rate with Claude Opus~4.6, compared with Claude~Code (58\%) and Cursor Composer~2 (61.7\%).  These results are notable because we did not tune our prompts or any model specifically for the Terminal Bench 2.0.
\end{abstract}

%----------------------------------------------------------------------
\section{Introduction}
\label{sec:intro}
%----------------------------------------------------------------------

Modern Large language models (LLMs), such as Anthropic's Claude Opus 4.7~\citep{anthropic2026opus47}, OpenAI's GPT 5.5~\citep{openai2026gpt55}, and Google's Gemini 3.1~\citep{google2026gemini31}, can generate code, reason about software architecture, and use developer tools~\citep{chen2021evaluating,roziere2023code}.  A growing body of work has explored how to harness these capabilities for autonomous software engineering, from single-session agents that resolve GitHub issues~\citep{yang2024sweagent,wang2024openhands} to industrial products marketed as AI software and general assistants ~\citep{copilot2021,cursor2024,devin2024,claudecode2025,codex2025,openclaw2025}. Yet using an LLM as a practical software engineering or general assistant still exposes several stubborn gaps: context windows are finite, a single mistake can derail an entire session, agents get stuck in dead ends, models generate AI slop, and generated changes are difficult to review or revert once applied to a live codebase.

We propose the \emph{KISS Agent Framework}, a stupidly simple AI agent framework containing around 2,500 lines of code for the core agent implementation.  We try to address the above-mentioned gaps through a structured system prompt (Section~\ref{sec:system-prompt}) and a five-layer agent hierarchy in which each layer solves exactly one concern:

\begin{enumerate}[nosep]
  \item \textbf{KISS Agent} --- budget-tracked
ReAct~\citep{yao2023react} loop with native function calling.
  \item \textbf{Relentless Agent} --- automatic summarization and continuation 
across sub-sessions.
  \item \textbf{Sorcar Agent} --- coding tools, browser automation,
and parallel sub-agent execution.
  \item \textbf{Chat Sorcar Agent} --- persistent multi-turn chat
sessions with history recall.
  \item \textbf{Worktree Sorcar Agent} --- git worktree isolation so
every task runs on its own branch.
\end{enumerate}

The name ``KISS'' reflects the \emph{Keep It Simple, Stupid} design philosophy in software engineering: each layer is small, each concern is isolated, and the overall system avoids unnecessary abstraction.

\begin{figure}[t]
\centering
\includegraphics[width=\textwidth]{KISS-Sorcar.png}
\caption{Screenshot of KISS Sorcar running as a VS~Code extension. The sidebar shows the agent's chat interface with real-time budget tracking, while the editor displays the code being modified.}
\label{fig:kiss-sorcar-screenshot}
\end{figure}

Table~\ref{tab:assistant-comparison} reproduces the project README's high-level comparison of KISS Sorcar with Claude Code and Cursor.

\begin{table}[H]
\centering
\caption{KISS Sorcar compared with Claude Code and Cursor, reproduced from the project README.}
\label{tab:assistant-comparison}
\scriptsize
\setlength{\tabcolsep}{3pt}
\renewcommand{\arraystretch}{1.15}
\begin{tabularx}{\textwidth}{>{\raggedright\arraybackslash}p{0.22\textwidth}>{\raggedright\arraybackslash}X>{\raggedright\arraybackslash}X>{\raggedright\arraybackslash}X}
\toprule
\textbf{Capability} & \textbf{KISS Sorcar} & \textbf{Claude Code} & \textbf{Cursor} \\
\midrule
\textbf{Interfaces} & CLI + VS Code extension + web/mobile app & CLI + mobile app & Custom VS Code \\
\textbf{AI Discovery} & \(\checkmark\) simply via prompt & \(\times\) & \(\times\) \\
\textbf{GEPA Prompt Optimization} & \(\checkmark\) simply via prompt & \(\times\) & \(\times\) \\
\textbf{Multiple models from multiple vendors in the same task} & \(\checkmark\) Mix OpenAI, Anthropic, Gemini, Together, Z.AI, Moonshot AI, OpenRouter, Claude Code CLI, and Codex CLI & \(\times\) Anthropic Claude models only & \(\times\) One model per task \\
\textbf{Primary focus} & \(\checkmark\) \textbf{Quality}---rigorous review, end-to-end tests & Speed and developer ergonomics & Speed \\
\textbf{Core Agents \# LoC} & \textbf{\textasciitilde{}2500} & Unknown & Unknown \\
\textbf{Models in bundled catalog} & 502 across 9 provider categories & Claude family only & Subset chosen by Cursor \\
\textbf{Bring your own API key / endpoint} & \(\checkmark\) Yes---keys stay on your machine & \(\checkmark\) Anthropic key & Routed through Cursor backend \\
\textbf{Open source} & \(\checkmark\) Apache-2.0 & \(\times\) Proprietary & \(\times\) Proprietary \\
\textbf{Price} & Free framework; pay only your chosen model provider & Subscription / API usage & Subscription \\
\textbf{Run on top of Claude Code / Codex CLI} & \(\checkmark\) \texttt{cc/*} and \texttt{codex/*} namespaces & N/A & \(\times\) \\
\textbf{Messaging and communication channels} & \(\checkmark\) 23 third-party agents, including Slack, Gmail, Phone Control, SMS, and WhatsApp & Partial: Slack, mobile Remote Control, and research-preview channels for Telegram, Discord, and iMessage; no documented built-in Gmail, WhatsApp, phone-call, or SMS channel & Partial: Slack and Microsoft Teams Cloud Agent integrations; no documented built-in Gmail, WhatsApp, phone-call, or SMS channel \\
\textbf{Terminal Bench 2.0 score} & \textbf{62.2\%} & 58\% & 61.7\% (Cursor agent) \\
\bottomrule
\end{tabularx}
\end{table}

We implement KISS Sorcar, a general-purpose assistant and an integrated development environment (IDE), on top of the KISS Agent Framework.  KISS Sorcar ships as three surfaces over the same daemon: a Visual Studio Code extension, a Claude-Code-style command-line interface (\texttt{sorcar}), and a browser/mobile web app---all running locally.  The CLI has both an interactive REPL mode (with slash commands such as \texttt{/help}, \texttt{/clear}, \texttt{/resume}, \texttt{/model}, \texttt{/cost}, \texttt{/skills}, \texttt{/mcp}, \texttt{/autocommit}, and file/folder \texttt{@}-mentions) and a one-shot non-interactive mode (\texttt{-t}/\texttt{-f}).  KISS Sorcar has browser support (using open-source Chromium and Playwright), multimodal support, Docker container support, OpenClaw-like features (whose discussion is beyond the scope of the paper), a mobile/web app, integrated MCP server management (\texttt{sorcar mcp add/list/auth/debug/\ldots}), and a library of Agent Skills loaded from \texttt{\textasciitilde{}/.kiss/skills} and project-local directories. KISS Sorcar is free and open-source (Apache-2.0), distributed on PyPI as \texttt{kiss-agent-framework}; all one needs is a model API key from a major LLM provider, such as Anthropic, OpenAI, Google, Together, MiniMax, or OpenRouter, or a custom local endpoint configured via \texttt{--endpoint}/\texttt{--header}.  The framework can also run on top of the Claude Code CLI (\texttt{cc/*} models) and the Codex CLI (\texttt{codex/*} models), and mix models from multiple vendors within a single task simply by issuing prompts.  We implemented this framework in roughly 4 months, and the repository is available at \url{https://github.com/ksenxx/kiss_ai}. The name ``Sorcar'' pays homage to P.~C.~Sorcar, the Bengali magician.  The engineering principles described in Section~\ref{sec:system-prompt} are encoded in the agent's system prompt.

KISS Sorcar has been built using itself.  The entire codebase---the KISS Agent framework, the Sorcar agent layers, the VS~Code extension, and the system prompt---was developed by KISS Sorcar operating on its own repository.  This self-hosting discipline provides a continuous-integration-style stress test: if the agent introduces a bug that impairs its ability to function, we ask the agent to fix it by analyzing the trajectory and code.  The simplicity of the layered architecture was both a prerequisite for and a consequence of this bootstrapping process. The simplicity of the framework reduced the number of bugs the agent introduced.  The five core agent classes remain compact: the KISS Agent comprises 423~lines, the Relentless Agent 412~lines, the Sorcar Agent 581~lines, the Chat Sorcar Agent 386~lines, and the Worktree Sorcar Agent 692~lines---a total of roughly 2,494~lines of code (counting only significant lines, i.e.\ excluding empty lines, comment-only lines, and docstrings of private methods).  Despite the Sorcar and Chat layers absorbing Docker-aware tool variants, parallel sub-agent orchestration, persistent chat-session bookkeeping, MCP integration, slash-command dispatch, and worktree concurrency safety, the framework has remained within this compact footprint and within the layered, single-concern design.

In the project, we deliberately prioritize output quality over speed.  In our experience, using a weaker or cheaper model often forces the developer to discard the agent's work and retry, ultimately increasing the total cost of completing a task.  Conversely, giving a frontier model adequate time to validate its own output---running linters, type checkers, and tests before declaring success---reduces the ``slop'' (low-quality, subtly incorrect code) common in faster but less thorough agents.  We expect token costs and inference latencies to continue to fall~\citep {gao2025trae}, making this quality-first posture increasingly practical.  In the meantime, the code produced by our agent is well-organized, simple, and idiomatic.

We evaluate on Terminal Bench~2.0 and achieve a 62.2\% overall pass rate using Claude Opus~4.6, compared with Claude~Code (58\%) and Cursor Composer~2 (61.7\%)~\citep{cursor2026composer2} on the same benchmark (Section~\ref{sec:evaluation}).  These results are notable because we did not tune our prompts or any model specifically for the Terminal Bench 2.0.

In KISS Sorcar, we deliberately kept the system simple.  We included only the agent technologies necessary for KISS Sorcar to function as a general-purpose software engineering assistant. We showed that a simple agent framework, without sophisticated agent technologies such as trajectory compaction and asynchronous multi-agent orchestration, was sufficient to build KISS Sorcar.  By building KISS Sorcar using itself and matching or exceeding both Cursor and Claude Code, we found that established software engineering techniques and principles are important for building reliable agent systems.

\paragraph{Outline.}
Section~\ref{sec:architecture} presents the five-layer agent architecture and its motivating design principles. Section~\ref{sec:evaluation} reports evaluation results on Terminal Bench~2.0.  Section~\ref{sec:system-prompt} details the system prompt. Section~\ref{sec:vscode-features} covers the user-facing features of the VS~Code extension, CLI, and web app. Section~\ref{sec:painless} illustrates painless software engineering through a real development session. Section~\ref{sec:related} discusses related work, and Section~\ref{sec:conclusion} concludes.

%----------------------------------------------------------------------
\section{Agent Architecture}
\label{sec:architecture}
%----------------------------------------------------------------------

We initially built the KISS Agent Framework to rapidly prototype and experiment with various prompt optimization techniques, such as Gepa~\citep{agrawal2025gepa}, and evolutionary algorithms for algorithmic and code optimization, such as AlphaEvolve~\citep{novikov2025alphaevolve} and OpenEvolve~\citep{openevolve2025}. We focused heavily on keeping the agent framework simple so we could rapidly prototype and experiment with ideas.  The simplicity of the framework also enabled coding agents to write simple, bug-free code.  We ultimately did not use any prompt optimization techniques when creating the system prompt for KISS Sorcar, as we hand-tuned it based on our long-term experience with KISS Sorcar and its behavior.  The KISS Agent Framework uses five agent layers in a layered architecture combining composition and inheritance. Each layer delegates upward for the concerns it does not own.

%......................................................................
\subsection{KISS Agent}
\label{sec:kiss-agent}
%......................................................................

The KISS Agent is the innermost execution unit.  It implements a standard ReAct loop~\citep{yao2023react} with the following characteristics.  

\begin{lstlisting}[style=kisspy,numbers=none,caption={A complete KISS agent with a single tool.},label={lst:simple-agent}]
from kiss.core.kiss_agent import KISSAgent

def calculate(expression: str) -> str:
    """Evaluate a math expression."""
    return str(eval(expression))

agent = KISSAgent(name="Math Buddy")
result = agent.run(
    model_name="gemini-2.5-flash",
    prompt_template="Calculate: {question}",
    arguments={"question": "What is 15% of 847?"},
    tools=[calculate]
)
print(result)  # 127.05
\end{lstlisting}

\textbf{Native function calling.}  We register tools as ordinary Python callables.  The agent builds an OpenAI-compatible tool schema once at setup time and caches it, avoiding redundant schema construction on every LLM call.  A special \texttt{finish} tool signals task completion and returns the result to the caller.

\textbf{Step, token, and budget tracking.}  At every step, the agent extracts input and output token counts from the API response, computes the dollar cost using a per-model pricing table, and updates both a local budget counter and a global (cross-agent) budget counter protected by a class-level lock.  The agent checks three limits before each step: the per-agent budget, the global budget, and the maximum step count.

\textbf{Error resilience.}  The agent retries transient API errors (rate limits, server errors) up to a configurable threshold of consecutive failures.  It detects non-retryable errors (authentication failures, permission denials) and raises them immediately.

\textbf{Non-agentic mode.}  When tools are not needed, the agent can run a single generation without the ReAct loop, which is useful for summarization or question-answering sub-tasks.

Listing~\ref{lst:simple-agent} shows a complete, working agent in under ten lines of code.  The developer defines an ordinary Python function (\texttt{calculate}), instantiates a \texttt{KISSAgent}, and calls its \texttt{run} method with a model name, a prompt template, template arguments, and a list of tools.  The framework automatically handles tool-schema generation, the ReAct loop, and budget tracking.

The KISS Agent is stateless across runs: each call to its run method resets the conversation, token counters, and tool registry.  This makes it safe to reuse a single agent instance for multiple sequential tasks.

%......................................................................
\subsection{Relentless Agent}
\label{sec:relentless-agent}
%......................................................................

The Relentless Agent wraps a KISS Agent in a continuation loop.  Its core contribution is the ability to execute tasks that exceed a single context window by breaking them into sub-sessions.

Rather than investing in context-compaction techniques, we adopt a simple continuation protocol: when a sub-session exhausts its context window or step budget, the agent produces a \emph{structured summary} of every action taken so far---chronologically ordered, with explanations and relevant code snippets---and a fresh sub-session resumes from that summary.  This approach is related in spirit to Reflexion~\citep{shinn2023reflexion}, which feeds verbal self-critiques back into subsequent trials, but uses Reflexion-like technique to continue a task.  While developing KISS Sorcar, we found in our experience that a na\"ive instruction to ``summarize the current context'' produced poor continuations; requiring a \emph{step-by-step chronological account with code snippets} improved coherence across sub-sessions. A potential limitation is that summaries may grow unwieldy for multi-day tasks; in practice, we have not encountered this problem even for tasks spanning several hours, but a thorough evaluation of summary scaling remains future work.

\textbf{Continuation protocol.}  The \texttt{finish} tool exposed to the inner KISS Agent accepts three fields: a success flag, a continue flag, and a summary.  When the agent sets \texttt{is\_continue=True}, the Relentless Agent starts a new sub-session with a fresh context window.  The prompt for the new session includes a chronologically ordered list of all prior attempt summaries and instructs the agent not to redo completed work.  The continuation prompt template is:

\begin{prompt}
# Task Progress (Continuation {continuation_number})

{progress_text}

# Continue
- Complete the rest of the task.
- **DON'T** redo completed work.
- If you have been retrying the same approach without progress,
  step back and rethink the strategy from scratch.
\end{prompt}

\textbf{Forced continuation on failure.}  If a sub-session raises an exception (e.g., the step limit is hit before the agent calls finish), the Relentless Agent does not abort.  Instead, it saves the full trajectory to a temporary file, spawns a separate summarizer agent to read the trajectory and produce a concise summary, and then uses that summary as the progress text for the next sub-session.  This ensures that even crashed sessions contribute useful context to subsequent attempts.  The summarizer receives the following prompt:

\begin{prompt}
# Summarizer

The trajectory of the agent is stored in the file: {trajectory_file}

# Instructions
- Read the trajectory file and analyze it.  The trajectory file
  could be large.
- Return a precise chronologically-ordered list of things the
  agent did with the reason for doing that along with relevant
  code snippets
\end{prompt}

\noindent To force the agent to self-continue before hitting the step limit, we augment the system prompt with an instruction that fires near the end of the budget:

\begin{prompt}
# MOST IMPORTANT INSTRUCTIONS
- **At step {step_threshold}: you MUST call
  finish(success=False, is_continue=True,
  summary="precise chronologically-ordered list of things
  the agent did with the reason for doing that along with
  relevant code snippets")** or if the task is not complete
  and you are at risk of running out of steps or context
  length.
- Work dir: {work_dir}
- Current process PID: {current_pid} -- NEVER kill this
  process.
\end{prompt}

%......................................................................
\subsection{Sorcar Agent}
\label{sec:sorcar-agent}
%......................................................................

The Sorcar Agent adds the tools that make the system useful for software development and general-purpose automation.

\textbf{Coding tools.}  We provide four core tools: a shell command executor with streaming output, a file reader, a precise string-based file editor, and a file writer.  The shell executor supports a configurable timeout, streams output to the user interface in real time, and respects a stop event that allows the user to cancel a running command.  Note that we kept the tool names (Bash, Edit, Write, Read) the same as in Claude Code because the underlying Anthropic models do not make mistakes with these tool names.

\textbf{Browser automation.}  A web-use tool provides programmatic browser control: navigating to URLs, reading page accessibility trees, clicking elements, typing text, pressing keys, scrolling, and taking screenshots.  This enables the agent to research documentation, verify deployed applications, and interact with web-based tools. It uses the open-source Chromium browser using the Playwright library. 

\textbf{Parallel sub-agents.}  An optional parallel execution tool spawns independent Sorcar Agent instances in a thread pool.  Each sub-agent gets its own LLM context and tool set.  This is useful for embarrassingly parallel tasks such as summarizing multiple files or researching independent topics.  We collect the results and return them in input order.  We do not activate this tool by default because in the IDE, we cannot stream multiple agent outputs coherently in the chat window.

\textbf{User interaction.}  An ask-user-question tool allows the agent to pause execution and request clarification from the user.  In the VS~Code integration, this renders as a text input in the sidebar; in CLI mode, it reads from standard input.

\textbf{Docker isolation.}  When a Docker image is specified, we replace the coding tools with Docker-aware variants that execute commands inside a container, providing an additional layer of sandboxing for untrusted tasks.

\textbf{Dynamic model switching with \texttt{set\_model}.}  A \texttt{set\_model} tool, registered automatically on every Sorcar Agent, lets the agent hand the live conversation off to a different LLM at any point during a task without restarting it.  The tool accepts a model name (for example \texttt{"gpt-5.5"}, \texttt{"claude-opus-4-7"}, or \texttt{"gemini-3.5-flash"}), constructs the new backend through the same factory that built the original model, copies the conversation history and the cumulative usage counters into the new backend, rebuilds the cached tools schema in the new provider's dialect (Anthropic and OpenAI schemas differ in subtle ways), persists the choice to disk so that the next task in the same chat session reuses it, and returns a confirmation string back to the agent loop.  If no model has been instantiated yet (deferred-startup case), it simply updates the agent's default \texttt{model\_name} for the next run; if the requested model is already active, it is a no-op.

This single tool is what enables \emph{multi-model, multi-vendor workflows inside a single task}.  Three idiomatic patterns recur in our usage: (i)~a \emph{scout-then-edit} pattern where a cheap fast model (e.g.\ \texttt{gemini-2.5-flash}) reads the repository, grep-searches for relevant call sites, and gathers context, then \texttt{set\_model("claude-opus-4-7")} hands off to a stronger reasoner for the actual edit; (ii)~a \emph{generate-then-review} pattern where a primary model produces a change and the agent calls \texttt{set\_model("gpt-5.5")} to switch to a second-opinion model that re-reads the diff, runs the tests, and verifies the change against the original request---this exact pattern is requested by the GEPA prompt-optimization task template in \texttt{src/kiss/SAMPLE\_TASKS.md} (``Use gpt-5.5 model (not codex) for thorough review of the work done at every step by the other model''); and (iii)~a \emph{cost-aware long-loop} pattern where an open-ended optimization or discovery task (Section~\ref{sec:long-running}) spends a cheap model on bookkeeping and journal-keeping but pays for a frontier model only at the decision points.  Because the conversation transcript and usage counters are carried across the switch, all subsequent budget accounting and continuation summaries remain coherent regardless of how many times the model is changed.

%......................................................................
\subsection{Chat Sorcar Agent}
\label{sec:chat-agent}
%......................................................................

The Chat Sorcar Agent adds multi-turn conversation persistence.

\textbf{Chat sessions.}  We assign each task to a chat session identified by a stable chat ID.  The agent persists tasks and their results to a local database (\texttt{sorcar.db}).  When a new task arrives within the same chat session, the agent loads prior tasks and results and prepends them to the prompt as numbered context entries, allowing the LLM to reference earlier work.

\textbf{Bounded chat context.}  To prevent unbounded growth of the prompt as a chat session accumulates many tasks, the agent caps the number of in-context history entries at \texttt{MAX\_TASKS=10}.  When the cap is exceeded, it preserves the first two entries (which typically establish the user's overall intent for the session) and the most recent entries, dropping the middle entries that are least likely to be referenced.  This keeps the context relevant and bounded while retaining both the session's framing and its current state.

\textbf{Session management.}  The agent supports three operations: starting a new chat (with a fresh chat ID), resuming a chat by task description (which looks up the corresponding chat ID), and resuming by explicit chat ID.  This enables both automatic session continuity in an IDE and manual session selection from the command line.

\textbf{Frequent task tracking.}  Each time a task is executed, the agent records the task description in a frequency table.  The IDE sidebar surfaces the most frequent tasks so users can re-issue them with one click, turning recurring requests (``run the test suite,'' ``regenerate the changelog'') into a click-to-replay experience.

\textbf{Metadata persistence.}  After each task, the agent records metadata including the model used, working directory, software version, token count, cost, and whether the task used parallel execution or worktree isolation.  This audit trail supports cost analysis and debugging.

%......................................................................
\subsection{Worktree Sorcar Agent}
\label{sec:worktree-agent}
%......................................................................

The Worktree Sorcar Agent is the outermost layer and the one that users interact with in the VS~Code extension and the Sorcar web app.  Its defining feature is git-worktree isolation.

\textbf{Branch-per-task.}  When a task starts, the agent creates a new git branch and a corresponding worktree directory.  The branch name encodes the chat ID and a timestamp for uniqueness.  All agent modifications happen inside the worktree; the user's main working tree remains untouched.

\textbf{Dirty-state preservation.}  If the user's main working tree has uncommitted changes, the agent copies them into the worktree and creates a baseline commit.  This ensures the agent sees the same state as the user, while keeping the user's actual index and working tree clean.  During merge, we use cherry-pick from the baseline commit to replay only the agent's changes, excluding the dirty-state snapshot.

\textbf{Concurrency safety.}  A per-repository file lock serializes the checkout, stash, merge, and pop sequence so that concurrent tabs in the IDE cannot interleave operations on the same repository. Thread-local storage isolates per-task state (stream parsing buffers, bash output buffers, recording state) so that stopping one task does not corrupt another.

\textbf{Crash recovery.}  We store all worktree state in git itself (branch names, git config entries) rather than in sidecar files.  On process restart, the agent queries git for any pending branch matching its chat ID prefix and reconstructs all instance attributes from git config, enabling recovery.

\textbf{Graceful fallback.}  If the working directory is not inside a git repository, if the repository has no commits, or if HEAD is detached, the agent falls back to direct execution without worktree isolation, ensuring it never fails due to git preconditions.


%----------------------------------------------------------------------
\section{AI Discovery, GEPA Optimization, and Repository Optimization}
\label{sec:long-running}
%----------------------------------------------------------------------

Beyond per-task coding assistance, KISS Sorcar can be driven as an autonomous \emph{research and optimization driver}: the user issues a single high-level objective with explicit numeric stopping conditions (a target accuracy, a target latency, a budget cap, a ``do not stop until\ldots'' clause), and the agent then runs an open-ended exploration loop---spawning experiments, reading logs, mutating its own prompts or code, journalling what worked and what failed---until the targets are met or the budget is exhausted.  The capabilities described in Section~\ref{sec:architecture} were not designed for this use case individually, but together they cover exactly what an open-ended driver needs: the Relentless Agent's structured-summary continuation (Section~\ref{sec:relentless-agent}) lets multi-hour or multi-day loops survive context exhaustion; the Sorcar Agent's parallel sub-agents and streamable shell executor (Section~\ref{sec:sorcar-agent}) parallelize the inner experiments and monitor long-running commands in real time; the \texttt{set\_model} tool lets a cheap model do the routine bookkeeping while a stronger model is engaged only at decision points; the Chat Sorcar Agent persists every tool call to \texttt{sorcar.db}, providing a replayable ground truth for self-reflection (Section~\ref{sec:chat-agent}); and the Worktree Sorcar Agent's branch-per-task isolation (Section~\ref{sec:worktree-agent}) keeps every trial reproducible and rollback-friendly.

We ship three sample task templates in \texttt{src/kiss/SAMPLE\_TASKS.md} that exercise this pattern.  Each is a single natural-language prompt that fits in a chat box; the agent does the rest.

%......................................................................
\subsection{AI Discovery}
\label{sec:ai-discovery}
%......................................................................

The AI Discovery template asks the agent to \emph{``discover the lightest and fastest AI model that will give the best accuracy and recall on the data at \texttt{<path>} at the cheapest price.''}  The prompt instructs the agent to (1)~analyze the data and search the internet extensively to propose the first few candidate models, (2)~split off 20\% of the data as an eval set that the discovery loop must never inspect, (3)~implement and benchmark each proposal on a smaller subset with fewer parameters first and then extrapolate, (4)~train on GPUs through a CLI such as Lambda Labs under a hard total budget (e.g.\ \$1000), and (5)~maintain a journal file of ideas already tried so that they are not repeated and so that successful ideas can be combined.  The prompt forbids memorization-style cheating in capitalized form---\emph{``MAKE SURE THAT YOU DO NOT DO REWARD HACKING OR CHEATING IN THE MODELS OR AGENTS YOU ARE IMPLEMENTING TO FIT DATA. YOUR SOLUTION MUST GENERALIZE BEYOND THE DATA PROVIDED.''}---and gives the agent a concrete stopping rule (e.g.\ 95\% recall on evals, fewer than 600~s per query, under \$50 per query amortized).  The driver ends by writing an HTML report with diagrams to \texttt{./reports/} and opening it in the user's default browser.

\textbf{Related work on AI-driven discovery.}  The AI Discovery template sits in a fast-moving line of work in which an LLM agent itself drives an outer optimization or search loop, edits source code or prompts as the unit of mutation, and verifies each candidate against an automated evaluator.  \emph{AlphaEvolve}~\citep{novikov2025alphaevolve}, the most prominent industrial instance, pairs a Gemini Flash/Gemini Pro ensemble with a programs database that selects parents under an evolutionary policy; deployed at Google it recovered $\sim$0.7\% of worldwide compute via a new Borg scheduling heuristic, sped up a Gemini matrix-multiplication kernel by 23\% and FlashAttention by up to 32.5\%, contributed a Verilog rewrite to an upcoming TPU, and discovered a 48-multiplication algorithm for 4$\times$4 complex matrix multiplication---the first improvement on Strassen's algorithm in this setting after 56 years.  \emph{OpenEvolve}~\citep{openevolve2025} is the open-source counterpart and adds MAP-Elites quality-diversity, island-based populations, an LLM ensemble with cascade evaluation, and an artifact side-channel that pipes stderr, profiler output, and LLM critique back into the next prompt.  \emph{ShinkaEvolve}~\citep{lange2025shinka}, from Sakana AI, attacks the sample-inefficiency that hampers earlier evolutionary coding agents through three additions---a parent-sampling rule that balances exploration and exploitation, code-novelty rejection sampling using both embedding similarity and an LLM-as-novelty-judge, and a bandit-based selection over an LLM ensemble; with these it sets a new state of the art for the 26-circle packing problem in only 150 samples, evolves an AIME-math agent scaffold in 75 generations, and discovers a Mixture-of-Experts load-balancing loss that outperforms DeepSeek's Global LBL in 30 generations.  \emph{ADRS} (AI-Driven Research for Systems)~\citep{cheng2025adrs}, introduced in the ``Barbarians at the Gate'' paper from UC Berkeley's Sky Computing Lab, argues that systems performance research is uniquely well-suited to this paradigm because verification reduces to running a simulator or real system against representative workloads; using OpenEvolve as a driver they report up to 5.0$\times$ runtime improvements and 50\% cost reductions over human-designed baselines on cross-region cloud scheduling, MoE inference, LLM-based SQL, and transaction scheduling, with most discoveries reached in under 12 hours at under \$20.  \emph{SkyDiscover}~\citep{liu2026skydiscover}, a follow-up framework from the same lab, factors the loop into a Context Builder, Solution Generator, Evaluator, and Solution Selector and instantiates two adaptive evolutionary strategies (AdaEvolve, EvoX); in $\sim$2.5K LoC it matches or exceeds AlphaEvolve on 6/8 math and 6/6 systems tasks and improves median scores on 172 Frontier-CS problems by $\sim$34\% over OpenEvolve, GEPA, and ShinkaEvolve.  \emph{GEPA}~\citep{agrawal2025gepa}, although phrased as a prompt optimizer, follows the same outer-loop pattern but uses natural-language reflection on full trajectories rather than scalar reward gradients: across six tasks it outperforms GRPO by up to 20\% while using up to 35$\times$ fewer rollouts, and outperforms MIPROv2 by over 10\%.  Finally, on the smallest end of the spectrum, Andrej Karpathy's \emph{AutoResearch}~\citep{karpathy2026autoresearch} reduces the loop to a single training file (\texttt{train.py}) edited by the agent, a single human-edited Markdown ``skill'' (\texttt{program.md}), and a fixed 5-minute training budget that makes architectural mutations directly comparable; this format demonstrates how cleanly an open-ended overnight ML-research loop can be expressed once an agent's coding tools and persistence story are robust.  Our AI Discovery template adopts the same fundamental contract as this family---natural-language objective, code-or-prompt as the search variable, automated evaluator as the verifier, journal of tried ideas, anti-reward-hacking and generalization clauses---but rather than ship a bespoke evolutionary harness, it relies on KISS Sorcar's general-purpose layered agent (Section~\ref{sec:architecture}) to provide every primitive the loop requires: budget-aware tool use, Relentless continuation across context windows, parallel sub-agents for independent candidate evaluation, branch-per-trial worktree isolation, replayable trajectories in \texttt{sorcar.db}, and dynamic \texttt{set\_model} switching between cheap rollout models and frontier reasoning models.  In this sense the AI Discovery template is a worked example of the ADRS thesis~\citep{cheng2025adrs} expressed inside a single general-purpose coding agent, with no separate orchestration framework.

This template exercises essentially every layer of the framework.  The Sorcar Agent's coding tools implement and benchmark each candidate; its web-use tool surveys arXiv, HuggingFace, and provider documentation for candidate models; its parallel sub-agents evaluate independent models on disjoint slices of the data.  The Relentless Agent's continuation protocol lets the journal of tried ideas and the train--eval--iterate loop persist across context windows.  The Worktree Sorcar Agent's branch-per-task isolation keeps each candidate model on its own branch so the user can review or roll back any individual trial.  The \texttt{set\_model} tool lets the agent run the bulk of the bookkeeping on a cheap model and only spend frontier-model budget on the experiment design and the result review steps where reasoning quality matters most.

%......................................................................
\subsection{GEPA Prompt Optimization}
\label{sec:gepa}
%......................................................................

The GEPA template instantiates the GEPA algorithm~\citep{agrawal2025gepa} to optimize the prompt used by a \texttt{ChatSorcarAgent} (Section~\ref{sec:chat-agent}).  The agent splits the labeled dataset 50/50 into a dev set and a val set and, on each outer iteration, samples a 100-example sub-val set (\texttt{sval}) and maintains a Pareto frontier of prompts under a \texttt{./pareto/} directory.  Each frontier node is itself a folder containing a \texttt{prompt.md} and a \texttt{score.json}; \texttt{score.json} is the list of \texttt{sval} session ids on which that prompt produces a correct prediction.  A new node is added only if its session set is not a subset of an existing node's, and nodes dominated by the new node are pruned afterward, so the frontier is monotone and bounded.

On each inner step the driver flips a coin: with probability 0.5 it picks a single frontier node, samples a 5-session minibatch from the dev set, runs the \texttt{ChatSorcarAgent} with that node's prompt, reads the trajectory of each failing session from \texttt{\textasciitilde/.kiss/sorcar.db}, and proposes a \emph{mutated} prompt that explicitly fixes the observed mistakes.  With the remaining probability it picks two frontier nodes and produces a \emph{crossover} prompt by merging the guidance from each.  In either case, any prompt that solves its minibatch is then evaluated on the full \texttt{sval} set and, if not dominated, joins the frontier.  The best-so-far prompt is mirrored to \texttt{./pareto/optimal.md} after every step, and the outer loop terminates after three consecutive iterations without improvement.

This template exposes the persistence layer as a first-class optimization signal: because every action of the \texttt{ChatSorcarAgent}---each tool call, each model response, each cost---is already written to \texttt{sorcar.db} by the Chat Sorcar Agent's metadata-persistence path (Section~\ref{sec:chat-agent}), the GEPA driver can replay any past trajectory and direct the mutator at the exact tool call that produced the failure, rather than guessing from a final answer alone.  The prompt also exploits \texttt{set\_model}: it requests a cheap model (e.g.\ Claude Opus~4.7) for the bulk of the rollouts and a frontier model (e.g.\ \texttt{gpt-5.5}, ``not codex'') only for the reflection and prompt-rewrite steps, and additionally for an end-of-step review by a different model than the one that produced the work.  As in AI Discovery, the prompt enforces the generalization clause---\emph{``YOUR SOLUTION MUST GENERALIZE BEYOND THE DATA PROVIDED''}---and the agent always appends ``Use internet search extensively at every step.'' to every candidate prompt so the optimized agent does not collapse into a closed-book solver.

%......................................................................
\subsection{Repository Optimization}
\label{sec:repo-opt}
%......................................................................

The repository optimization template targets an existing code base together with a runnable command (a build, a benchmark, a server, or a training script): \emph{``Can you run the command \texttt{<command>} in the background and monitor its output in real time to optimize the code at \texttt{<folder\_or\_url>} with respect to the following metrics: speed, accuracy, recall, cost?''}  The agent is explicitly allowed to add diagnostic code that emits finer-grained metric values, but the prompt also requires it to remove that diagnostic code after the optimization is complete, so the final commit is clean.  As in AI Discovery, the agent maintains a journal of ideas tried, repeated failures, and successful directions to combine; the user supplies concrete numeric targets and the agent does not stop until they are met.

This pattern exercises three KISS Sorcar capabilities together that no other layer can substitute for.  First, the Sorcar Agent's shell executor with streaming output and a cancellable stop event lets the agent launch a long-running command, watch its output in real time, and abort it the moment a code change invalidates the current run---without this, every speculative optimization would have to wait for a full benchmark to finish.  Second, the Worktree Sorcar Agent's branch-per-task isolation means each candidate optimization lands on its own branch, so a regression discovered three iterations later can be reverted with a single \texttt{git} operation while preserving progress made on later, unrelated branches.  Third, the \texttt{set\_model} tool lets the agent profile and edit on a cost-efficient model and switch to a stronger model only for the harder algorithmic rewrites and for a final review pass.

%......................................................................
\subsection*{Common pattern}

All three templates share the same skeleton: a high-level natural-language objective, explicit numeric stopping conditions, a journal file of tried ideas, parallel exploration where possible, branch-per-trial isolation, an anti-reward-hacking clause, an explicit instruction to use internet search extensively at every step, and a final HTML report written into \texttt{./reports/} and opened in the user's default browser.  The agent's system prompt (Section~\ref{sec:system-prompt}) already encodes the engineering disciplines---read-before-modify, lint-and-test before \texttt{finish}, no fabricated source counts, no shortcuts---so the user need only state the objective, the targets, and any external constraint (such as a GPU CLI to use).  In our experience, this open-ended-driver use case is where the layered architecture pays off most: each layer's narrow concern---budget, continuation, parallelism, persistence, isolation, dynamic model selection---is exactly what a long-running self-directed exploration requires, and the simplicity of the framework leaves no room for the driver to silently corrupt its own state over a multi-hour run.


%----------------------------------------------------------------------
\section{Evaluation on Terminal Bench~2.0}
\label{sec:evaluation}
%----------------------------------------------------------------------

Before we discuss the system prompt in a lengthy section, we describe the evaluation outcome of KISS Sorcar on the Terminal Bench 2.0, which was also recently used by the Cursor agent of Composer 2.0.

We evaluate our system on Terminal Bench~2.0,\footnote{\url{https://www.tbench.ai/}} a benchmark comprising 89 diverse terminal-based programming tasks, ranging from building legacy compilers and configuring servers to solving cryptanalysis challenges and training machine-learning models.  Each task runs in an isolated Docker container; a separate verifier automatically judges the result.  We use the Harbor\footnote{\url{https://github.com/harbor-framework/harbor}} framework to orchestrate execution, and Claude Opus~4.6 as the underlying LLM.  We do not modify the general system prompt or inject Terminal Bench 2.0-specific instructions during the evaluation.  We carried out our evaluation on a 2025 MacBook Air 15" with an M4 processor and 24GB RAM.

\subsection{Setup}

We run 5 independent trials per task.  The agent is \texttt{SorcarHarborAgent}, a thin Harbor adapter that installs and invokes the Sorcar CLI inside each container.  We hard-skip 9 tasks that we verified to be infeasible for Opus~4.6 across 6+ prior attempts (e.g.\ CompCert compilation, Windows~3.11 GUI installation, video OCR) to save time and token cost.  Skipped tasks still count as failures.

\subsection{Aggregate Results}

Table~\ref{tab:tb2-aggregate} summarizes the aggregate statistics.

\begin{table}[h]
\centering
\caption{Terminal Bench~2.0 aggregate results (89 tasks, 5 trials each, Claude Opus~4.6).}
\label{tab:tb2-aggregate}
\begin{tabular}{lr}
\toprule
\textbf{Metric} & \textbf{Value} \\
\midrule
Total tasks              & 89 \\
Overall pass rate        & 62.2\% (277/445) \\
pass@any (at least 1/5 passes) & 78.7\% (70/89) \\
pass@all (all 5 pass)    & 43.8\% (39/89) \\
Always-fail tasks        & 19 \\
Always-pass tasks        & 39 \\
Mixed-result tasks       & 31 \\
Median cost per trial    & \$0.45 \\
Mean cost per trial      & \$0.90 \\
Median duration per trial & 202\,s \\
Mean duration per trial  & 446\,s \\
\bottomrule
\end{tabular}
\end{table}

The 62.2\% overall pass rate is comparable to other agents using the same underlying model: at the time of writing, Claude~Code (also Opus~4.6) scores approximately 58\% on the Terminal Bench~2.0 leaderboard, and Cursor's Composer~2---a custom fine-tuned model trained with large-scale reinforcement learning~\citep{cursor2026composer2}---achieves 61.7\%.  This suggests that the layered architecture and the structured system prompt described in Sections~\ref{sec:architecture}--\ref{sec:system-prompt} contribute beyond what the base model alone provides.

\subsection{Task-Level Breakdown}

\textbf{Consistently solved tasks (39 of 89).} These include cryptanalysis (FEAL differential), game-playing (chess best move), git operations (leak recovery), server configuration (gRPC key-value store, PyPI server, NGINX logging), data processing (resharding), formal verification (Coq~\texttt{plus\_comm}), ML inference (HuggingFace model serving, LLM batching scheduler), and system emulation (QEMU startup).  These tasks span systems, security, data engineering, and formal methods.

\textbf{Consistently failed tasks (19 of 89).} The failures cluster into three categories: (1)~tasks requiring graphical or multimedia capabilities unavailable in the container (video processing, Windows~3.11 GUI, MTEB leaderboard scraping, extracting moves from video), (2)~tasks demanding very long or resource-intensive builds that exceed the container's time or memory limits (CompCert, Doom for MIPS, Caffe CIFAR-10, training fastText on Yelp data), and (3)~tasks with niche domain-specific requirements that the model struggles to satisfy (DNA insertion, OCaml GC patching, polyglot C/Python binaries, protein assembly, cell segmentation).

\textbf{Mixed-result tasks (31 of 89).} Tasks such as \texttt{write-compressor} (3/5), \texttt{crack-7z-hash} (4/5), and \texttt{feal-linear-cryptanalysis} (4/5) succeed in most trials but occasionally fail due to non-determinism in the model's reasoning or timing-sensitive environment interactions.  Conversely, \texttt{cancel-async-tasks} (1/5) and \texttt{dna-assembly} (1/5) succeed rarely, suggesting they are at the boundary of the model's capability.

\textbf{Leaderboard context.}  KISS Sorcar does not score as high as other coding agents reported at the Terminal Bench 2.0 leaderboard, but these results are notable because we did not tune our prompts or any model specifically for the Terminal Bench 2.0.  We used the general system prompt and Claude Opus 4.6 without modification.  Regarding the lower score compared to other coding agents, recent analysis has found widespread cheating on popular agent benchmarks, including Terminal Bench~2.0: the top three submissions commit harness-level cheating (e.g.\ leaking verifier code or answer keys into the agent's environment), and task-level cheating (e.g.\ Googling answers, mining git history, hardcoding test outputs) affects 28+ submissions across 9 benchmarks~\citep{stein2026cheatingblog,stein2026detecting}. Separately, we discovered, using an automated benchmark audit agent, that 45 confirmed hacking solutions across 13 widely used benchmarks exhibited process-isolation failures, answer leakage, and weak test assertions that allow perfect scores without solving a single problem~\citep{wang2026trustworthyblog}.

%----------------------------------------------------------------------
\section{The System Prompt}
\label{sec:system-prompt}
%----------------------------------------------------------------------

The system prompt is a structured document that governs the agent's behavior across all tasks.  It is not a generic instruction to ``be helpful'' but a specification of engineering practices.

\textbf{XML-tagged structure.}  The prompt is organized using XML tags that delimit each concern: \texttt{<identity>}, \texttt{<visibility\_constraint>}, \texttt{<tool\_rules>}, \texttt{<web\_research>}, \texttt{<code\_style>}, \texttt{<workflow>}, \texttt{<testing>}, \texttt{<pre\_finish\_verification>}, and \texttt{<sorcar\_specific>}.  All three major LLM providers---Anthropic~\citep{anthropic2025prompting}, OpenAI~\citep{openai2025prompting}, and Google~\citep{google2025prompting}---recommend XML tags for structuring complex prompts: they create unambiguous section boundaries that models parse as structural markers rather than content, reducing the chance that the model misinterprets an instruction from one section as applying to another.

\textbf{Identity-first ordering.}  We place the agent's identity and visibility constraints before any behavioral rules.  This follows the recommended prompt ordering for frontier models: the model should know \emph{what it is} before learning \emph{what to do}.

We describe the key sections below.

%......................................................................
\subsection{Identity and Visibility}
%......................................................................

The prompt opens with two XML-tagged sections that establish who the agent is and how it communicates with the user:

\begin{prompt}
<identity>
You are KISS Sorcar, an AI General Assistant and IDE
developed by Koushik Sen (ksen@berkeley.edu).
Repo: https://github.com/ksenxx/kiss_ai
Version: 2026.6.28

Your sole goal is completing the user's task accurately
and thoroughly.  Be rigorous, check facts, and produce
high-quality work.
</identity>

<visibility_constraint>
The user cannot see your thoughts, reasoning, scratchpad,
intermediate tool outputs, or assistant prose. The ONLY
thing the user sees is the string you pass to
finish(summary=...).  Compose the full detailed answer
directly inside the summary string of finish(). When
answering informational questions, include the complete
answer in the summary, not a meta-description of what
was done.

**Bad** (meta-description): "Greeted the user and asked
what they'd like to work on. Awaiting a specific task."
**Good** (actual content): "Hi! I'm KISS Sorcar, ready
to help. What would you like to work on?"

The summary must contain the actual content the user
should see, not a third-person narration of what
happened.
</visibility_constraint>
\end{prompt}

\noindent Each section addresses a distinct concern:

\textbf{Identity placement.}  The \texttt{<identity>} block appears first in the prompt, before any behavioral rules.  This follows the recommended prompt ordering for frontier models~\citep{anthropic2025prompting,openai2025prompting}: the model should know \emph{what it is} before learning \emph{what to do}.  The identity block also consolidates directives that were previously scattered as aggressive imperatives (``BE RELENTLESS,'' ``BE RIGOROUS,'' ``CHECK FACTS,'' ``NO AI SLOP'') into a single calm sentence: ``Be rigorous, check facts, and produce high-quality work.''  Research from all three major providers indicates that positive, explanatory framing is more reliable than capitalized commands with frontier models.

\textbf{Task focus.}  ``Your sole goal is completing the user's task accurately and thoroughly'' anchors the model on the task at hand and discourages meta-commentary, tangential exploration, and unsolicited clarification questions that consume tokens without making progress.  It also instructs the model to treat errors as obstacles to overcome rather than reasons to stop.

\textbf{Visibility constraint as a separate section.}  The \texttt{<visibility\_constraint>} block is separated from tool rules because it governs a different concern: not \emph{how} to use tools, but \emph{what the user can see}.  Without this instruction, the model may ``tell'' the user something in an intermediate message and then assume the user has seen it, leading to confusion when the user asks for information the model believes it already provided.  The clause ``not a meta-description of what was done'' prevents the model from returning vague summaries like ``Fixed the bug in Y'' instead of showing the actual fix; it forces the model to include concrete details, results, and outputs in the summary.  A second audit found that 3~out of 91~production tasks still produced meta-descriptions (e.g., ``Greeted the user and asked what they'd like to work on'' instead of the actual greeting).  The current version now includes explicit \textbf{Bad}/\textbf{Good} examples directly in the prompt---contrasting a third-person meta-description against the actual greeting the user should see---to make the distinction concrete and unambiguous, and closes with the directive that ``the summary must contain the actual content the user should see, not a third-person narration of what happened.''

%......................................................................
\subsection{Tool Rules}
%......................................................................

Tool usage rules are explicit and mechanical:

\begin{prompt}
<tool_rules>
## Tool Usage

- Use Write() for new files; Edit() for small changes.
- Run Bash synchronously with timeout_seconds (default
  120s). On timeout, retry with a higher value.  For
  commands exceeding 10 minutes, run in background,
  redirect output to a file, and poll periodically.
- Use go_to_url() for browser navigation.
- Read large files in chunks.
- **Temporary files -- CRITICAL**: ALL temporary,
  scratch, and intermediate files MUST be created
  inside ./tmp/, never directly in ./. This includes
  research notes, file-information dumps, downloaded
  artifacts, build outputs, and any other transient
  file. Create ./tmp/ if it doesn't exist.  Before
  calling finish(), delete every temporary file you
  created in ./tmp/ (but not the directory itself if
  it was pre-existing).
- When multiple independent tool calls are needed, make
  them all in the same turn to maximize parallelism.
  When calls depend on prior results, sequence them
  across turns.

## Context and Continuation

- If running out of context or steps, do not rush. Call
  finish(is_continue=True) to pause and resume the task
  in a new context.
</tool_rules>

- If there is ambiguity or under specification in the
  user task, search the internet to find the most
  reliable and modern solution to resolve the ambiguity.
\end{prompt}

Each tool rule addresses a specific failure mode.  An earlier revision of this section opened with an explicit definition---``\textbf{PWD denotes current working directory} and does not refer to a directory named PWD''---added after we observed the model creating a literal \texttt{PWD/} subdirectory inside the workspace.  The current revision drops the disambiguation entirely by replacing every occurrence of \texttt{PWD/} in the prompt with the shell-conventional relative form \texttt{./}~(e.g., \texttt{./tmp/}, \texttt{./SORCAR.md}), which frontier models interpret unambiguously as the working directory.  This is a small example of a recurring simplification pattern: when a prompt rule exists solely to disambiguate a confusing notation, replacing the notation is preferable to explaining it.

\textbf{``Use Write() for new files; Edit() for small changes.''}  Without this distinction, the model may use \texttt{Write()} to overwrite an existing file with a slightly modified version, losing content it forgot to include.  By reserving \texttt{Write()} for new files and requiring \texttt{Edit()} for modifications, the instruction ensures that changes are surgical and that unchanged portions of a file are never at risk.

\textbf{Bash timeout guidance.}  LLMs frequently launch shell commands without considering their runtime.  A compilation or test suite that takes five minutes will time out at the default 30-second shell timeout in most agent frameworks, causing spurious failures.  The instruction to use 120~seconds as the default, retry with higher timeouts on timeouts, and run long-running commands in the background with output redirected to a file provides a mechanical protocol that handles common cases without requiring the model to estimate runtime from first principles.

\textbf{``Use go\_to\_url() for browser navigation.''}  The agent has access to multiple tools that could plausibly interact with the web (shell-based \texttt{curl}, a Python program, the browser tool, etc.).  This clause eliminates ambiguity by specifying which tool to use for browser-based interactions.

\textbf{``Read large files in chunks.''}  Reading a 10,000-line file in a single tool call consumes a large fraction of the context window.  By instructing the model to read files in chunks, the prompt prevents context window exhaustion caused by a single-file read, preserving capacity for the rest of the task.

\textbf{Temporary files (CRITICAL).}  A separate, emphatically marked clause requires that \emph{all} temporary, scratch, and intermediate files---research notes, file-information dumps, downloaded artifacts, build outputs---live inside \texttt{./tmp/} (the \texttt{tmp/} subdirectory of the working directory) and never directly in the working directory.  Without this directive, the model creates temporary files in unpredictable locations (the system \texttt{/tmp}, the home directory, or scattered throughout the project) or sometimes directly in the project root, polluting the working tree with artifacts that are difficult to distinguish from legitimate project files.  Centralizing temporary files in a known directory makes cleanup mechanical and predictable.  The companion clause requires the agent to create \texttt{./tmp/} if it does not exist and to delete every temporary file it created before calling \texttt{finish}, while preserving the \texttt{./tmp/} directory itself if it was pre-existing---a deliberately narrow cleanup contract that prevents the agent from accidentally deleting unrelated files that happened to be in the same directory.

\textbf{Ambiguity protocol.}  A trailing instruction outside the \texttt{<tool\_rules>} block---``If there is ambiguity or under specification in the user task, search the internet to find the most reliable and modern solution to resolve the ambiguity''---channels the model's natural tendency to confabulate when underspecified into the structured web-research workflow described below.  Rather than letting the model guess (the default behavior) or stop and ask the user (which interrupts the workflow), this directive turns ambiguity into a research task that produces an auditable artifact.

\textbf{Parallel tool calls.}  The instruction to ``make them all in the same turn'' when tool calls are independent reduces latency by allowing the framework to execute multiple tool calls concurrently.  Without this instruction, the model tends to issue one tool call per turn even when the calls are independent, wasting round trips.

\textbf{Context and continuation.}  When the context window is nearly full, LLMs exhibit a ``rush to finish'' behavior: they skip verification steps, make hasty edits, and call \texttt{finish} with an incomplete result.  The continuation instructions redirect that urgency into the continuation protocol (Section~\ref{sec:relentless-agent}), ensuring that a clean handoff to a new sub-session produces better results than a frantic attempt to squeeze everything into the remaining tokens.

%......................................................................
\subsection{Pre-flight Checks}
%......................................................................

We begin every task with mandatory reads, then enforce a read-before-modify discipline:

\begin{prompt}
## Mandatory First Actions for project related
tasks -- CRITICAL

**Your VERY FIRST tool call** in every task MUST be
Read("./SORCAR.md") and follow the instructions
in SORCAR.md with highest priority.

## Pre-flight Checks

**Read before modify rule -- NON-NEGOTIABLE**: You MUST
call Read(file_path) on every file BEFORE calling
Edit(file_path) on it. Never Edit a file you have not
Read in the current session.

**Use the file tools, never shell substitutes --
CRITICAL**: To VIEW the contents of any file, you MUST
use the Read() tool. It is FORBIDDEN to inspect or dump
file contents through Bash using cat, sed -n, head,
tail, awk, more, less, nl, grep over a whole file,
echo "$(<file)", or a for-loop of cat. These do NOT
satisfy the read-before-modify rule and waste context
-- always call Read() instead (use max_lines/chunking
for big files). To MODIFY any file, you MUST use the
Edit() or Write() tools. It is FORBIDDEN to edit files
in place through Bash using sed -i, perl -i,
awk ... > file, tee, or output redirection (>, >>) onto
a source/tracked file. Bash may still be used for
non-file-content operations (running tests, ls, grep -l
to find files, git, builds, moving/removing files).

Read relevant source files when the task depends on
existing architecture. If referenced files, commands, or
config don't exist, stop and ask the user rather than
guessing.

**When fixing bugs, issues, or race conditions: write an
end-to-end test that reproduces the problem first, then
fix the code, then verify the test passes.**
\end{prompt}

\noindent The ``Mandatory First Actions'' section ensures that project-specific overrides are loaded before any work begins.  This section has undergone several rounds of empirical refinement.  In an earlier version that merely mentioned the read inside a separate ``Self-Improvement Loop'' section, the agent read \texttt{SORCAR.md} as its first tool call only 46\% of the time.  After promoting the read to a dedicated section with imperative language, compliance improved in integration tests but regressed in production: analysis of 91~real tasks showed that 92\% never read \texttt{SORCAR.md}.  Two evasion patterns emerged: (1)~the agent used \texttt{Bash("cat ...")} instead of the \texttt{Read()} tool, and (2)~the agent combined the file read with other commands in a single \texttt{Bash} call.  Explicit prohibitions and ordinal position mandates raised compliance to 100\% in integration tests.

The current version collapses the mandatory first action to a single read of \texttt{SORCAR.md}, accompanied by the rider ``follow the instructions in SORCAR.md with highest priority.''  This is a deliberate simplification: \texttt{SORCAR.md} contains project-specific instructions that can override or extend the general system prompt, and loading it first lets the rest of the task be interpreted under the active project's rules.  Earlier revisions of the prompt also mandated a second read of a per-project \texttt{USER\_PREFS.md} preferences file paired with a ``Self-Improvement Loop'' section that asked the agent to update the file at task end; both have since been removed because the self-learning store accumulated stale project facts (Section~\ref{sec:self-improvement-prompt}) that drifted out of sync with the evolving code, and the read-only half of the protocol added prompt bulk without measurable benefit.  The prohibitions against \texttt{Bash}-based reading have since been promoted into an explicit, enumerated rule of their own---``Use the file tools, never shell substitutes''---in the Pre-flight Checks section.

Each pre-flight check targets a specific category of avoidable error:

\textbf{``Read before modify rule.''}  The most common source of agent-introduced bugs is modifying a file based on an incorrect assumption about its current contents.  The model may ``remember'' an older version of the file from its training data, or it may extrapolate from a partial reading.  By requiring a fresh read immediately before any edit, the instruction ensures that the model operates on the file's current state rather than a stale mental model.  The companion clause ``Read relevant sources if the task depends on existing architecture'' extends this rule from individual files to architectural context.  A task like ``add a caching layer to the database module'' requires understanding not just the file to be modified, but also how callers interact with it, what interfaces it provides, and what invariants it assumes.  The instruction prevents the model from jumping straight to code generation without understanding the broader context.

\textbf{``Use the file tools, never shell substitutes.''}  The read-before-modify discipline is only effective if the model actually routes file access through the \texttt{Read()}, \texttt{Edit()}, and \texttt{Write()} tools, where the framework can track, render, and (for edits) verify the operation.  In practice, the model repeatedly evaded this by reaching for the shell instead---dumping a file with \texttt{cat}, \texttt{sed -n}, \texttt{head}, \texttt{tail}, \texttt{awk}, \texttt{nl}, or a \texttt{grep} over the whole file, and editing in place with \texttt{sed -i}, \texttt{perl -i}, \texttt{tee}, or output redirection (\texttt{>}, \texttt{>>}).  These shell substitutes bypass the read tracking entirely (so a subsequent \texttt{Edit()} fires against a file the framework never saw being read) and waste context by spilling raw file contents into the transcript.  The current prompt therefore enumerates the forbidden commands explicitly for both viewing and modifying, clarifies that editing a file viewed only through a forbidden shell command is a ``double violation'' (the file must be \texttt{Read()} first, then \texttt{Edit()}/\texttt{Write()}), and---crucially---carves out the legitimate uses of Bash that must remain available (running tests, \texttt{ls}, \texttt{grep -l} to locate files, \texttt{git}, builds, and moving or removing files).  Enumerating both the prohibited and the permitted shell uses prevents the model from over-applying the rule and abandoning Bash for tasks where it is the correct tool.

\textbf{``If referenced files, commands, or config don't exist, stop and ask the user rather than guessing.''}  LLMs have a strong tendency to confabulate: when asked to modify a file that does not exist, the model will often proceed as if it does, producing edits against phantom content.  This instruction converts a silent failure (incorrect edits applied to a nonexistent file, which silently creates it) into an explicit clarification request.

\textbf{``Write an end-to-end test that reproduces the problem first, then fix the code, then verify the test passes.''}  This instruction mandates a test-first discipline for bug fixes, specifying the three-step sequence explicitly.  The motivation is two-fold: first, a test that reproduces the bug provides concrete verification that the fix is correct (the test should pass after the fix and fail before it).  Second, writing the test forces the model to understand the bug precisely before attempting a fix, reducing the risk of an ad~hoc patch that addresses a symptom rather than the root cause.  Specifying ``end-to-end test'' rather than just ``test'' reinforces the no-mocks testing discipline described in Section~\ref{sec:system-prompt}.

%......................................................................
\subsection{Code Style Guidelines}
%......................................................................

The prompt encodes a minimalist code philosophy:

\begin{prompt}
## Code Style

Write simple, clean, readable code with minimal
indirection. These rules exist because over-abstracted
code is harder to debug and maintain.

- Organize code across multiple files grouped by
  functionality.
- Prefer named functions, classes, and module-level
  helpers over closures and lambdas. Closures obscure
  control flow; use explicit parameter passing instead.
- Eliminate unnecessary attributes, locals, config vars,
  tight coupling, and attribute redirections.
- Eliminate redundant abstractions and duplicate code.
- Public methods must have full docstrings.
- Fix root causes, not symptoms. Before writing code,
  ask: is this simple, elegant, general, and minimal?
- Write documentation only when the task explicitly
  requires it.
\end{prompt}

\noindent Each guideline addresses a specific anti-pattern commonly
exhibited by LLM-generated code:

\textbf{``These rules exist because over-abstracted code is harder to debug and maintain.''}  This rationale sentence is deliberate.  Anthropic's prompting guide recommends providing motivation behind instructions to help the model generalize correctly~\citep{anthropic2025prompting}.  When the model understands \emph{why} a rule exists, it can apply the underlying principle to edge cases that the rule does not explicitly cover.

\textbf{``Write simple, clean, readable code with minimal indirection.\ Organize code across multiple files grouped by functionality.''}  LLMs tend to over-engineer solutions, introducing unnecessary abstractions, helper classes, and levels of indirection.  Simple code is easier to review, test, and maintain.  LLMs also often pile new code onto whichever file is currently being edited, producing 2,000-line modules that conflate unrelated concerns.  This directive nudges the model toward a modular layout in which each file has a single, coherent responsibility.

\textbf{``Prefer named functions, classes, and module-level helpers over closures and lambdas.\ Closures obscure control flow; use explicit parameter passing instead.''}  LLMs reach for closures whenever a small piece of state needs to be carried alongside a function---producing nested \texttt{def}s that capture mutable variables from the enclosing scope.  Such closures are difficult to test in isolation, opaque to type checkers, and a frequent source of subtle bugs due to the late binding of captured variables.  Rather than a blanket prohibition (as in the earlier version of this prompt), the revised instruction uses positive framing: it tells the model what to \emph{prefer} and explains \emph{why}, which is more effective with frontier models that interpret instructions literally.  The instruction steers the model toward explicit data structures (plain functions with arguments, classes with attributes), which are easier to reason about, easier to test, and play well with our no-mocks testing discipline.

\textbf{``Eliminate unnecessary attributes, locals, config vars, tight coupling, and attribute redirections.''}  LLMs frequently introduce intermediate variables that serve no purpose---for example, assigning a return value to a local variable only to immediately return it on the next line, or storing a constant in a configuration file when it is used in exactly one place.  When the model adds a feature that touches multiple files, it may introduce imports, shared global state, or cross-module function calls that create tight coupling.  An attribute redirection occurs when an object stores a reference to another object solely to forward method calls to it---for example, \texttt{self.x = other.x} at construction time, creating two paths to the same value.  This single consolidated rule addresses all of these anti-patterns.

\textbf{``Eliminate redundant abstractions and duplicate code.''}  LLMs sometimes create utility functions or classes that duplicate existing functionality; this instruction reminds the model to check for existing implementations before creating new ones.

\textbf{``Public methods must have full docstrings.''}  While the prompt generally discourages unnecessary documentation (see the last item), public methods are the API surface that other developers and modules depend on.  Documentation on public methods is not optional---it specifies the contract.

\textbf{``Fix root causes, not symptoms.\ Before writing code, ask: is this simple, elegant, general, and minimal?''}  LLMs frequently apply symptom-level fixes: adding a null check where the real problem is that a variable should never be null, or catching an exception where the real problem is that the caller passes invalid arguments.  This instruction forces the model to trace the causal chain to the root and fix it there.  The companion metacognitive instruction asks the model to pause and evaluate its plan before committing to an implementation, spending more inference-time compute on design and reducing the likelihood of producing an unnecessarily complex first draft.

\textbf{``Write documentation only when the task explicitly requires it.''}  Claude Opus 4.6 tends to generate many documentation files.  This instruction prevents the behavior.

%......................................................................
\subsection{Deep Work Rules}
%......................................................................

\begin{prompt}
## Deep Work

- For tasks involving "align", "match", or "make
  consistent": read the target state fully before
  editing. Never edit based on vague recollection.
- Use concrete values, not indirections. Read file Y
  first, then write the specific values into file X.
- List concrete planned changes before executing
  multi-part work.
- Every meaningful change needs a concrete verification
  method (test, grep, CLI check).
\end{prompt}

\noindent The deep work rules address a failure mode where the model interprets an instruction loosely and makes changes that are directionally correct but concretely wrong:

\textbf{``For `align'/`match'/`make consistent': read the target state before editing.''}  When a user says ``make file~A consistent with file~B,'' the model often reads file~A, infers what file~B probably contains, and edits A based on that inference---without ever reading B.  This instruction mandates reading the target first, ensuring that the alignment is based on concrete facts rather than assumptions.

\textbf{``Use concrete values, not indirections (read Y first, then write specific values into X).''}  A related failure mode occurs when the model's plan says ``update X to match Y'' but the model never resolves what Y actually is.  The instruction requires the model to first read~Y, extract the specific values, and then write those values into~X.  This eliminates a class of errors where the model's mental model of~Y differs from reality.

\textbf{``List concrete planned changes before executing multi-part work.''}  When a task requires changes to multiple files, executing them one at a time without a plan leads to inconsistencies: the model may change a function signature in one file but forget to update a caller in another.  Listing all planned changes before executing any of them forces the model to consider the full scope of the change and identify dependencies.

\textbf{``Every meaningful change needs a concrete verification method.''}  A change without a verification method is a change that cannot be confirmed to work.  This instruction requires the model to pair each change with a specific check---a test, a grep for the expected pattern, a CLI command that exercises the changed behavior---ensuring that the change can be validated programmatically rather than by visual inspection of a diff.

%......................................................................
\subsection{Planning for Complex Tasks}
%......................................................................

The planning instructions use a complexity threshold---three or more files, cross-module changes, or architectural work---to decide when formal planning is required:

\begin{prompt}
## Complex Task Planning

For work spanning 3+ files, crossing module boundaries,
or changing architecture:

1. List every file to change and why.
2. State the exact intended change per file.
3. Identify dependencies and execution order.
4. State the verification method per change.

Skip this planning step for simple single-file
modifications.
\end{prompt}

\noindent Each planning step targets a specific failure mode:

\textbf{``List files to change and why.''}  This forces the model to enumerate the full blast radius of a change before touching any file.  Without this step, the model often discovers mid-task that additional files need changes, leading to incomplete or inconsistent modifications.

\textbf{``State exact intended change per file.''} Listing files alone is insufficient; the model must also articulate \emph{what} will change in each file.  This converts a vague plan (``update the database module'') into a concrete specification (``add a \texttt{cache\_ttl} parameter to \texttt{DatabaseClient.\_\_init\_\_}, modify the query method to check the cache before hitting the database, add a cache invalidation method).

\textbf{``Identify dependencies and execution order.''}  Some changes must precede others: a new utility function must be written before callers can import it, a migration must run before code that depends on the new schema.  Identifying these dependencies prevents the model from applying changes in an order that produces intermediate states where the code does not compile, or tests do not pass.

\textbf{``State verification method per change.''}  The verification requirement from the Deep Work section is reinforced here at the planning stage, ensuring that verification is planned alongside the changes rather than treated as an afterthought.

\noindent The escape clause---``Skip for simple single-file tasks''---avoids the overhead of planning trivial changes.  Requiring a formal plan for a one-line typo fix would waste tokens and slow down the agent without any compensating benefit.

%......................................................................
\subsection{Testing Instructions}
%......................................................................

The testing section is perhaps the most opinionated:

\begin{prompt}
## Testing

- Run lint and typecheckers; fix all errors including
  pre-existing ones.
- Aim for 100% branch coverage on new and modified code.
- Write end-to-end tests only. Do not use mocks,
  patches, fakes, or test doubles. Each test must be
  independent and verify actual behavior.
- **DO NOT** write structural tests which assert on
  the source code.
- After modifications, run only the impacted tests.
- To confirm race conditions: add a random sleep (<0.1s)
  before the suspected racing statements.
- **CRITICAL**: Before running tests, count the number
  of tests. If the number of tests is more than 100,
  split the set of tests equally by the number of test
  methods into number of cores - 2 and run all splits
  in parallel using run_parallel tool.
\end{prompt}

\noindent Each testing instruction addresses a specific concern:

\textbf{``Run lint and typecheckers; fix all errors including pre-existing ones.''}  Before committing any change, the agent must ensure it does not introduce lint violations or type errors.  This catches a broad class of issues---unused imports, type mismatches, style violations---that would otherwise accumulate across tasks.  The clause ``including pre-existing ones'' prevents the model from rationalizing existing errors as ``not my problem'' and calling \texttt{finish} with a passing result despite a broken build.  The instruction makes the agent responsible for the entire codebase health, not just the delta it introduced.

\textbf{``Aim for 100\% branch coverage on new and modified code.''}  LLMs tend to write happy-path tests that cover the main code path but ignore error handling, edge cases, and early-return branches.  The 100\% target forces the model to write tests for every branch, including error paths and boundary conditions.  Moreover, such tests help with regression---developers can use AI coding agents with less risk that changes will break existing program behavior.  The wording ``aim for'' rather than ``achieve'' acknowledges that perfect coverage is not always feasible, while still setting an ambitious target.

\textbf{``Write end-to-end tests only.\ Do not use mocks, patches, fakes, or test doubles.''}  This is the most opinionated rule.  Mock tests that verify code calls certain methods in a certain order test the implementation, not the behavior.  A test suite built on mocks can pass with flying colors while the system is fundamentally broken, because the mocks hide the real dependencies.  End-to-end tests that exercise actual behavior are more expensive to run but provide stronger evidence that the system works.  Moreover, writing end-to-end tests forces the model to reason about the system's actual dependencies, often enabling the agent to find additional bugs.  The distinction between unit and end-to-end tests matters: a unit test in isolation may verify that a function produces the right output for a given input, but an end-to-end test verifies that the function works correctly within the larger system---with real file I/O, real database connections, and real inter-module interactions.

\textbf{``Each test independent, verifying actual behavior.''}  Test independence means that running tests in any order produces the same results.  Tests that depend on shared state or execution order are brittle and difficult to debug when they fail.  ``Verifying actual behavior'' reiterates that tests should assert on observable outcomes (return values, side effects, system state) rather than implementation details.

\textbf{``Only run impacted tests after modifications.''}  Running the full test suite after every small change is wasteful when only a few modules are affected.  For a large project, a full test run may take minutes, and doing it after every edit adds up to significant wasted time and compute.  This instruction directs the model to identify which tests are affected by its changes and run only those, improving iteration speed.

\textbf{``To confirm races: add random sleep (<0.1s) before racing statements.''}  Race conditions are notoriously difficult to reproduce because they depend on precise timing.  By inserting small sleep delays at strategic points, the model can widen the race window and make the bug manifest deterministically during testing.  The 0.1-second upper bound keeps the test fast while still being sufficient to expose most races.

\textbf{``Do not write structural tests which assert on the source code.''}  LLMs frequently fall back on \emph{structural} assertions when they cannot easily exercise a behavior: asserting that a particular function exists, that a class has a certain attribute, that an import appears in a specific order, or that a file contains a specific substring.  Such tests do not exercise the system---they merely encode the current shape of the code.  They fail spuriously after harmless refactors (renaming a private helper, reordering imports, inlining a function) while providing no evidence that the program actually works.  Worse, they create a false sense of coverage: a green suite that consists mostly of structural assertions can coexist with a completely broken runtime.  The explicit prohibition steers the agent back to behavioral assertions that survive refactoring and genuinely guard correctness.

\textbf{Parallel test execution above 100 tests (CRITICAL).}  Once a suite grows beyond roughly a hundred tests, running it serially dominates the agent's iteration time and inflates token cost (because the model waits, then re-reads, then re-reasons about a long output).  The instruction tells the agent to count tests first, and if the count exceeds 100, to partition the test set evenly into \emph{number of cores~\(-\)~2} shards and dispatch them concurrently via the \texttt{run\_parallel} tool.  The ``cores~\(-\)~2'' reservation deliberately leaves headroom for the agent's own process and for the IDE so that aggressive parallelism does not starve the foreground experience.  The 100-test threshold is the empirical break-even point at which the orchestration overhead of \texttt{run\_parallel} (spawning shards, collecting and merging results) is amortized by the wall-clock savings; below the threshold, sequential execution is simpler and faster overall.  This rule is marked CRITICAL because, like the lint/typecheck obligation, agents otherwise revert to the path of least resistance---a single \texttt{pytest} invocation---and pay a large hidden cost on every long-running task.

%......................................................................
\subsection{Self-Improvement via the Framework}
\label{sec:self-improvement-prompt}
%......................................................................

Earlier revisions of \texttt{SYSTEM.md} contained an explicit \emph{Self-Improvement Loop} section that instructed the agent, before calling \texttt{finish}, to update a \texttt{USER\_PREFS.md} file with any reusable user preference, project convention, file location, or other durable fact discovered during the task.  The section also encoded the curation rules: no code snippets or symbol names in entries, skip one-off task details, remove conflicting older entries, and keep the file small because its content was loaded into every subsequent task's context.  The motivation was that the preferences file would let the agent accumulate project knowledge across sessions even though each session starts with a fresh context window.

The current \texttt{SYSTEM.md} no longer contains a \texttt{\#\# Self-Improvement Loop} section and the \texttt{USER\_PREFS.md} file has been removed from the project entirely.  The mechanism was removed because the self-learning store accumulated \emph{stale} information about the project.  Facts that were true when an entry was written---file locations, API endpoints, project conventions, and design decisions---silently went out of date as the codebase evolved, and because the stored entries were injected into every subsequent task's context, later tasks acted on outdated assumptions.  In several cases this caused the agent to reintroduce bugs that had already been fixed: it ``learned'' a workaround or invariant, that fact later became false, and the agent kept applying the obsolete knowledge.  The preferences file thus accumulated noise faster than useful invariants and actively regressed the code rather than improving it.  Project-specific instructions turned out to be better captured in the per-repository \texttt{SORCAR.md} override file, which the user controls directly and keeps current.  Removing the self-learning mechanism entirely shrinks the prompt, eliminates a brittle cross-session state file that drifted out of sync with the code, and concentrates project knowledge in a single human-curated place that cannot silently go stale.

%......................................................................
\subsection{Pre-Finish Verification}
%......................................................................

Before declaring a task complete, the agent must pass a structured verification checklist:

\begin{prompt}
## Pre-Finish Verification -- CRITICAL

Before calling finish(success=True):

1. Re-read and verify every modified file.
2. If you created or modified ANY .py, .ts, .js, .css,
   .tsx, or .jsx file in this session: you MUST run
   uv run check --full and fix all errors. This is not
   optional. Do NOT call finish without running this
   command first. If the project doesn't use uv, run
   the equivalent lint/typecheck command.
3. Check each user requirement against what was
   delivered.
4. **Clean up temporary files -- MANDATORY**: You MUST
   delete every temporary file you created in ./tmp/
   during this session (research notes, information-*.md,
   file-information-*.md, scratch scripts, downloaded
   artifacts, etc.).  Explicitly run Bash("rm -f
   ./tmp/<each-file-you-created>") and then
   Bash("ls ./tmp") to confirm they are gone. Do NOT
   call finish(success=True) while any temp file you
   created still remains. Do NOT delete files you did
   not create.
5. If any check fails, keep working.
6. After 3 failed retries of the same fix approach,
   step back and rethink from scratch.
\end{prompt}

\noindent Each step in this checklist addresses a specific way agents declare premature success:

\textbf{``Re-read and verify every modified file.''}  This is the analog of a code review performed by the agent on its own work.  The model may have introduced a typo, forgotten to close a bracket, or made an edit that looked correct in the diff but was wrong in the full-file context.  Re-reading the file after all edits are complete catches these errors.

\textbf{``If you created or modified ANY code file: you MUST run uv run check.''}  The original instruction (``Run required checks; fix failures'') was too vague: analysis of 91~real tasks showed that 70\% of code-modification tasks (7/10) skipped lint/typecheck entirely.  The revised version enumerates the exact file extensions that trigger the obligation and names the exact command to run, converting a soft guideline into an unambiguous, verifiable mandate.  A trailing fallback clause---``If the project doesn't use uv, run the equivalent lint/typecheck command''---generalizes the rule to projects with different toolchains while keeping \texttt{uv run check -{}-full} as the canonical instruction.

\textbf{``Check each user requirement against delivery.''}  The model may have completed a task that it \emph{thinks} satisfies the user's request, but actually misses a requirement.  This instruction forces a systematic comparison between the original task description and the delivered result, catching gaps and misinterpretations.

\textbf{``Clean up temporary files.''}  The temporary-files directive in the Tool Usage section requires the agent to write scratch files into \texttt{./tmp/}, but without an explicit cleanup step at the end of the task, those artifacts persist and pollute the project's working tree over time.  This step mandates deletion of every temporary file the agent created during the session.  The companion clause ``Do NOT delete files you did not create'' is essential: the \texttt{tmp/} directory may contain artifacts from other sessions or from the developer's own scratch work, and an indiscriminate \texttt{rm -rf tmp/*} would destroy unrelated content.  Together, the two clauses form a narrow, audit-friendly cleanup contract.  The current revision additionally embeds an explicit command template---\texttt{Bash("rm -f ./tmp/<files-you-created>")}---directly inside the bullet, so the agent has a concrete, copy-pasteable cleanup command and does not have to invent its own (occasionally over-broad) deletion strategy.  It further requires the agent to run \texttt{Bash("ls ./tmp")} afterward to confirm the files are gone, and forbids calling \texttt{finish(success=True)} while any agent-created temporary file remains---turning the cleanup from an easily skipped suggestion into a verified post-condition.

\textbf{``If any check fails, keep working.''}  Without this instruction, the model may call \texttt{finish(success=True)} even when it knows a check has failed, rationalizing that the failure is ``minor'' or ``unrelated.''  The instruction makes the rule absolute: no finishing until all checks pass.

\textbf{``After 3 failed retries of same fix, rethink from scratch.''}  LLMs can enter repetitive loops where they apply the same incorrect fix repeatedly, each time hoping for a different result.  The three-retry threshold forces the model to break out of such loops by abandoning the current approach and reconsidering the problem from first principles.  This is analogous to the debugging heuristic ``if you've been staring at the same code for twenty minutes, you're looking in the wrong place.''

%......................................................................
\subsection{Web Research Protocol}
%......................................................................

When the agent needs external knowledge, the prompt prescribes a structured research workflow rather than allowing ad-hoc browsing:

\begin{prompt}
## Web Research

When a task requires searching the internet, researching
a topic, or answering questions that benefit from current
information:

- Visit at least 10 distinct websites per research
  session. Do not stop early or rationalize visiting
  fewer. **This is a hard requirement -- you MUST
  visit 10 sites, not 5 or 10.**
- You MUST use go_to_url() to visit each site. Do NOT
  use Bash("curl ...") or Bash("wget ...") as a
  substitute for visiting websites. Using curl/wget to
  fetch pages does not count toward the 10-site
  requirement.
- Procedure:
  1. Create ./tmp/information-{unique_id}.md with
     header: # Web Research -- Websites visited: 0/10
  2. Per site visited: (a) use go_to_url() to visit the
     site, (b) extract information, (c) use Edit() to
     append ## [N/10] URL + extracted information to the
     file, (d) use Edit() to update the header counter.
  3. Do not proceed to synthesis until the counter
     reaches 10.
  4. If results dry up, try different queries, synonyms,
     official docs, GitHub repos/issues, Stack Overflow,
     blogs, Reddit, papers, and API references.
  5. After reaching 10, review all findings and
     synthesize.
- Ask the user for login help when a page requires
  authentication.

This requirement applies to research and
information-gathering tasks. For pure code edits, bug
fixes, or file modifications where you already have
sufficient context, proceed directly.
\end{prompt}

\noindent The rationale is a two-phase separation between \emph{collection} and \emph{synthesis}.  LLMs tend to anchor on the first few results they encounter, which biases their solutions toward a narrow slice of the design space.  By forcing the agent to accumulate a broad set of information into a file \emph{before} reasoning about it, the protocol counteracts anchoring bias and encourages the model to consider diverse approaches.  An earlier version of this prompt required visiting at least 30 websites, but the threshold was lowered to 10 after we observed that the 30-site requirement frequently inflated token cost and wall-clock time without proportionally improving answer quality on routine research tasks; 10 sites remain enough to span official documentation, primary sources, secondary commentary, and a few divergent perspectives.  The resulting information file serves as an auditable artifact of what the agent considered.  The structured procedure with a counter header (\texttt{\# Web Research --- Websites visited: 0/10}) and per-site entries (\texttt{\#\# [N/10] URL}) addresses an empirically observed failure mode in which the model claims to have ``visited many sites'' after only a handful of fetches; a concrete counter forces the model to verify the actual number visited before declaring the collection phase complete.  An additional emphatic clause---``This is a hard requirement -- you MUST visit 10 sites, not 5 or 10''---was added after we observed the model rationalizing partial compliance (``I visited 7 sites, which is close enough'').  The instruction to try ``different queries, synonyms, official docs'' when results dry up prevents the agent from giving up prematurely on a narrow set of search terms.

An additional evasion pattern observed in production was the agent using \texttt{Bash("curl ...")} to fetch web pages instead of \texttt{go\_to\_url()}, thereby bypassing the browser-based visit counter while still gathering information.  Analysis of 91~real tasks found 4~research tasks that visited fewer than the required minimum number of sites (ranging from 1--5~URLs), with one task (1179) using \texttt{curl} to fetch content from 15+~sites without triggering \texttt{go\_to\_url}.  The explicit prohibition of \texttt{curl}/\texttt{wget} for research closes this loophole.

A second audit uncovered two additional failure modes.  First, the mandatory information file (\texttt{tmp/information-\{id\}.md}) was created in 0~out of 4~research tasks---the agent skipped file creation entirely and answered from memory or from a handful of sites.  The current version adds explicit mandatory language: ``The file is your proof of work---if it doesn't exist when you call finish, you violated this rule.''  Second, one task claimed ``a synthesis of 30+~sources'' in its result summary while having visited only 4~URLs---a fabricated source count.  A new ``Real-Time Data'' section addresses both hallucination (answering current weather, news, or stock questions from stale training data---one task reported news from the wrong year) and source fabrication (``Do NOT fabricate or exaggerate source counts'').

The login instruction addresses a practical obstacle in web research: many websites require authentication before revealing their content.  Rather than silently skipping gated pages or hallucinating their contents, we instruct the agent to ask the user for help with login.

\textbf{Scoped applicability.}  The final paragraph---``This requirement applies to research and information-gathering tasks.  For pure code edits, bug fixes, or file modifications where you already have sufficient context, proceed directly''---is an important addition.  Without this exemption, the agent would perform 10 website visits even for simple one-line code fixes where the necessary context is already in the file being edited, wasting tokens and tool-call budget.  The exemption allows the agent to skip research when the task is purely mechanical, while still enforcing thorough research for tasks that benefit from external information.

%......................................................................
\subsection{File Browsing Protocol}
%......................................................................

When a task requires understanding multiple source files before making changes, the prompt prescribes the same two-phase collect-then-synthesize discipline used for web research, but applied to the local file system:

\begin{prompt}
## File Browsing

When exploring unfamiliar code, collect information and
code snippets in ./tmp/file-information-{unique_id}.md
as you go relevant for the task, then review the
collected material and think deeply before acting.
\end{prompt}

\noindent This instruction addresses a failure mode distinct from the web research case.  When an agent must read many project files to understand a codebase before making changes, it tends to read a file, form a hypothesis, and immediately begin editing---anchoring on the first few files it encounters and missing relevant context in files it never opens.  Worse, each file read consumes context window tokens; by the time the agent has read enough files to understand the full picture, it may have already spent most of its context window on the raw file contents, leaving little room for reasoning and code generation.

The file browsing protocol counteracts both problems.  By writing a structured summary of each file's relevant information into a temporary markdown file, the agent externalizes its understanding into a compact artifact that persists across context boundaries.  The instruction to collect ``without overthinking'' is deliberate: during the collection phase, the agent should extract and record facts (function signatures, class hierarchies, call sites, invariants) rather than analyze or plan.  Analysis happens in the second phase, when the agent reads its own summary file and reasons about the collected information as a whole.

This two-phase separation provides three benefits.  First, it prevents premature commitment: the agent cannot start editing until it has surveyed the relevant files, reducing the risk of changes that are locally correct but globally inconsistent.  Second, the summary file is typically much smaller than the raw source files, freeing up context window capacity for subsequent reasoning and editing phases.  Third, the summary file serves as an auditable artifact---the developer can inspect it to verify that the agent considered the right files and extracted the right information before making changes.

%......................................................................
\subsection{Desktop Application Control}
%......................................................................

The agent can interact with graphical desktop applications using screenshots, keyboard, and mouse:

\begin{prompt}
## Desktop Apps

Interact with desktop applications using screenshots, keyboard, and mouse. Do not launch VS Code or its extensions.
\end{prompt}

\noindent This instruction enables the agent to operate GUI applications (Preview, browsers, graphical diff tools) when command-line alternatives are insufficient.  The explicit prohibition on launching VS~Code prevents a recursive loop: since the agent runs \emph{inside} a VS~Code extension, launching another VS~Code instance or modifying extension state from within the agent could corrupt the host session or create deadlocks.  Note that modern LLMs support desktop control abilities, and we are merely exploiting them.

%......................................................................
\subsection{Sorcar-Specific Overrides}
%......................................................................

A final section provides project-specific instructions that are injected when the agent operates on its own codebase:

\begin{prompt}
## Sorcar-specific

- Lint/typecheck/format: `uv run check --full`. Tests:
  `uv run pytest -v` (timeout 900s).
- The injectable instructions are available at
  ~/.vscode/extensions/ksenxx.kiss-sorcar-2026.6.28/
    kiss_project/src/kiss/INJECTIONS.md
- Your SYSTEM.md (the system prompt) is located at
  ~/.vscode/extensions/ksenxx.kiss-sorcar-2026.6.28/
    kiss_project/src/kiss/SYSTEM.md
- KISS Sorcar paper:
  https://github.com/ksenxx/kiss_ai/blob/main/
    papers/kisssorcar/kiss_sorcar.tex
- Third-party agents: kiss/agents/third_party_agents
- Claude SKILLS: kiss/agents/claude_skills. You can
  use them as necessary.
- **If you create any artifact that the user can use
  after the task is over, you MUST create them in a
  directory and add the directory contents to git.**
- MAINTAIN a PROGRESS.md across agent sessions logging
  details of all the steps you have done so far from
  the start with explanation and relevant code snippets.
  Clear the contents of PROGRESS.md when a new task
  (not subtask) begins.
- **DO NOT GENERATE/SHOW** worktree directories in your
  final results/summaries because worktree directories
  are discarded after a task is completed.  Rather show
  the directories relative to the main repo.
- Authenticate unauthenticated third-party agents; ask
  the user only when a page requires human
  authentication.
\end{prompt}

These overrides serve nine purposes.  First, they specify the exact toolchain commands for the KISS project itself (\texttt{uv run check -{}-full}, \texttt{uv run pytest -v} with a 900-second timeout), eliminating guesswork about which linter, formatter, or test runner to use.  Second, the agent is told the on-disk location of an \emph{injectable instructions} file---\texttt{\textasciitilde/.vscode/extensions/ksenxx.kiss-sorcar-\$VERSION/kiss\_project/src/kiss/INJECTIONS.md}---which the framework may concatenate onto the system prompt at runtime to deliver tricks and workarounds discovered after the extension shipped without requiring a new release; pointing the agent at this file lets it inspect or update those injected directives directly.  Third, the agent is told the on-disk location of its own \texttt{SYSTEM.md} (under the same extension directory, ending in \texttt{/SYSTEM.md}) so that questions about its own prompt or behavior can be answered by reading the canonical file rather than from memory.  This pointer is paired with a URL to the paper's \LaTeX{} source, giving the agent two complementary self-references: the system prompt for operational rules and the paper for design rationale.  Fourth, the instructions expose a third-party agent integration layer: the agent is told where third-party agents live (\texttt{kiss/agents/third\_party\_agents}).  When a third-party agent requires authentication, the agent handles it autonomously and only prompts the user when a page genuinely requires human credentials---reducing unnecessary interruptions while maintaining security.  Fifth, Claude SKILLS~\citep{claudecode2025} are bundled in \texttt{kiss/agents/claude\_skills} with the note ``You can use them as necessary,'' giving the agent access to Anthropic's skill library for common software engineering patterns; the permissive framing (rather than a mandate) lets the agent select skills opportunistically when a task matches one.  Sixth, the \texttt{SORCAR.md} override mechanism---now enforced by the Mandatory First Actions section---allows per-repository instructions to further customize the agent's behavior, forming a hierarchy: general system prompt $\rightarrow$ Sorcar-specific instructions $\rightarrow$ repository-specific \texttt{SORCAR.md}.  This approach also enables us to not hardcode \texttt{SORCAR.md} in the KISS Sorcar code, removing one more dependency.  In the \texttt{SORCAR.md}, one can include more markdown files for further instructions.  Seventh, a directory-with-git-commit rule requires that any artifact the user can use after the task ends be created inside a directory and added to git, ensuring that durable outputs are version-controlled and easy to locate rather than scattered as loose files.  Eighth, the agent is required to maintain a \texttt{PROGRESS.md} across sub-sessions logging every step taken so far with explanations and relevant code snippets---and to clear it when a new top-level task (not a subtask) begins; this acts as a persistent scratchpad that bridges Relentless Agent continuations and lets a fresh sub-session pick up exactly where the previous one left off without re-deriving prior decisions.  Ninth, the worktree-presentation rule (``DO NOT GENERATE/SHOW worktree directories in your final results/summaries... rather show the directories relative to the main repo'') prevents the agent from quoting transient worktree paths that are discarded after merge; user-visible references must point at paths in the main repository so that copy-pasted file names and shell commands remain valid after the task completes.

%----------------------------------------------------------------------
\section{User-Facing Features}
\label{sec:vscode-features}
%----------------------------------------------------------------------

We release our system as three coordinated surfaces over a single local daemon: a VS~Code extension, a Claude-Code-style command-line interface (\texttt{sorcar}), and a browser/mobile web app.  While the underlying agent architecture (Sections~\ref{sec:architecture} and~\ref{sec:system-prompt}) already differs from existing AI coding assistants, the user-facing design introduces several features that differ from those in existing IDE assistants such as GitHub Copilot~\citep{copilot2021}, Cursor~\citep{cursor2024}, Windsurf~\citep{windsurf2024}, Devin~\citep{devin2024}, and Aider~\citep{aider2023}.  We describe these features below.

%......................................................................
\subsection{Real-Time Budget Accountability}
\label{sec:budget-feature}
%......................................................................

AI coding assistants typically operate on a subscription model (Copilot, Cursor) or a per-seat pricing model (Devin, Windsurf), both of which obscure the per-task cost.  The developer has no visibility into how many tokens a task consumed or how much it cost.

Our extension displays real-time cost tracking in the sidebar: input tokens, output tokens, cache hits, dollar cost, and elapsed time are updated at every agent step.  We enforce both per-task and global budget ceilings.  If a task exceeds its budget, the agent raises a hard error rather than silently accumulating charges.  This transparency allows developers to make informed decisions about which tasks to delegate to the AI and how to structure prompts for cost efficiency.  The KISS Agent also appends the current usage in the context so that the model is fully aware of its limits.

%......................................................................
\subsection{Integrated Browser Automation}
\label{sec:browser-feature}
%......................................................................

Our extension includes a browser automation tool that allows the agent to navigate to URLs, read accessibility trees, and click elements, type text, press keys, scroll, and take screenshots--all controlled programmatically from within a VS~Code task via Playwright. We render a live browser preview in a Chromium browser, allowing the developer to watch the agent interact with web applications in real time.

This capability enables use cases that most IDE assistants do not support: verifying a deployed web application after a code change, filling out web forms as part of a testing workflow, scraping documentation to inform a code generation task, or interacting with web-based developer tools (CI dashboards, issue trackers) without leaving the editor.

%......................................................................
\subsection{Interactive CLI with Slash Commands and \texttt{@}-Mentions}
\label{sec:cli-feature}
%......................................................................

The \texttt{sorcar} command-line interface runs in two modes.  In \emph{interactive} mode (no \texttt{-t}/\texttt{-f} flag), it presents a Claude-Code-style REPL that connects as a thin terminal client to the same local \texttt{sorcar web} daemon used by the VS~Code extension.  In \emph{non-interactive} mode (\texttt{-t} or \texttt{-f} supplied), it runs a single task and exits.  Either mode honors flags for model selection (\texttt{-m}), custom endpoints and headers (\texttt{-e}/\texttt{--header}), per-task budget caps (\texttt{-b}), working-directory pinning (\texttt{-w}), worktree isolation (\texttt{--worktree}/\texttt{--no-worktree}), automatic commit on task finish (\texttt{--auto-commit}/\texttt{--no-auto-commit}), and toggles for browser tools (\texttt{--no-web}) and parallel sub-agents (\texttt{--no-parallel}).

The interactive REPL adds two affordances borrowed from modern coding assistants but rare in open-source agents.  First, file and folder \emph{\texttt{@}-mentions} with ranked project-file completion let the user paste path-aware context into the prompt without leaving the keyboard.  Second, a small set of slash commands provides direct control over the agent: \texttt{/help} lists every command, \texttt{/clear} (alias \texttt{/new}) starts a fresh chat, \texttt{/resume} reopens a prior chat by id, \texttt{/model} (and \texttt{/model list}) switches or enumerates models mid-session, \texttt{/cost} (aliases \texttt{/usage}, \texttt{/context}) prints the running token and dollar totals plus context-window utilization, \texttt{/skills} and \texttt{/mcp} introspect loaded Agent Skills and configured MCP servers, \texttt{/autocommit} toggles the worktree's auto-commit policy, \texttt{/commands} lists user-defined Markdown slash commands, and \texttt{/exit} (alias \texttt{/quit}) terminates the session.  Custom Markdown slash commands are auto-loaded from \texttt{\textasciitilde/.kiss/commands}, \texttt{<project>/.kiss/commands}, and the corresponding Claude directories, so users can register reusable prompt templates per-project or globally without modifying the Sorcar source.

%......................................................................
\subsection{Extensibility: MCP, Skills, and Third-Party Agents}
\label{sec:ext-feature}
%......................................................................

KISS Sorcar exposes three pluggable extension points that let users add capabilities without forking the framework.

\textbf{Model-Context-Protocol (MCP) servers.}  The CLI's \texttt{sorcar mcp} subcommand registers, lists, inspects, and authenticates MCP servers in either user scope (\texttt{\textasciitilde/.kiss/mcp.json}) or project scope (\texttt{<project>/.kiss/mcp.json}, with backward compatibility for \texttt{<project>/.mcp.json}).  Servers may use stdio, HTTP, or SSE transports.  An OAuth~2.1 flow (\texttt{sorcar mcp auth}) implements dynamic client registration with PKCE and persists tokens under \texttt{\textasciitilde/.kiss/mcp\_auth/}; a \texttt{sorcar mcp debug} command dumps a server's capabilities, tools (with input schemas and granted permissions), resources, and prompts.  Discovered tools become first-class \texttt{KISSAgent} tools alongside the built-in Bash/Edit/Read/Write tools.

\textbf{Agent Skills.}  At startup the framework scans \texttt{\textasciitilde/.kiss/skills}, \texttt{<project>/.kiss/skills}, Anthropic Claude skill directories, \texttt{.agents/skills}, and the bundled Sorcar skills directory.  Each skill is exposed to the agent as a callable that the model may invoke when the task description matches the skill's natural-language trigger.  This mirrors Anthropic's Claude SKILLS~\citep{claudecode2025} convention so that the same skill files work in both ecosystems.

\textbf{Third-party messaging and automation agents.}  KISS Sorcar ships with 23~third-party agents under \texttt{src/kiss/agents/third\_party\_agents} that wrap real-world communication and consumer-product surfaces---BlueBubbles, Discord, Feishu, Gmail, Google Chat, iMessage, IRC, LINE, Matrix, Mattermost, Microsoft Teams, Nextcloud Talk, Nostr, Phone Control, Signal, Slack, SMS, Synology Chat, Telegram, Tlon, Twitch, WhatsApp, and Zalo.  It additionally ships a Govee smart-home CLI for controlling IoT lights (on/off, brightness, color, and color temperature) via the Govee Developer API.  Each third-party agent declares the credentials or OAuth flow it needs; when missing, the agent attempts to authenticate autonomously and only escalates to the user when a step genuinely requires a human (e.g., entering a two-factor code).  This breadth distinguishes KISS Sorcar from Claude~Code (whose documented channels include Slack, mobile remote control, and research-preview Telegram/Discord/iMessage but no built-in Gmail, WhatsApp, phone-call, or SMS) and from Cursor (whose Cloud Agent integrations cover Slack and Microsoft Teams).

\textbf{VS~Code surfaces for runtime customization.}  Two additional plain-text files customize the VS~Code extension without rebuilding it.  \texttt{\textasciitilde/.kiss/INJECTIONS.md} (seeded from the bundled \texttt{src/kiss/INJECTIONS.md}) supplies entries for the sidebar's ``Tricks'' dropdown---one entry per \texttt{\#\# Trick} section---so that workarounds discovered after the extension shipped can be injected into the system prompt at task launch.  \texttt{\textasciitilde/.kiss/SAMPLE\_TASKS.md} (seeded from the bundled \texttt{src/kiss/agents/vscode/SAMPLE\_TASKS.md}) supplies the welcome-screen sample-task chips---one chip per \texttt{\#\# Task} section.  Deleting either file regenerates it from the bundled defaults on next launch, so customization is non-destructive.

%......................................................................
\subsection{Multi-Model, Multi-Vendor Workflows}
\label{sec:multi-model-feature}
%......................................................................

KISS Sorcar ships with a catalog of 501 models spanning eight provider categories---70 OpenAI, 13 Anthropic, 23 Gemini/Google, 77 Together AI, 5 MiniMax, 303 OpenRouter, 3 Claude~Code CLI (exposed under the \texttt{cc/*} namespace), and 7 OpenAI Codex CLI (\texttt{codex/*} namespace) entries---of which 485 are generation-capable, 321 are function-calling-capable, and 7 are embedding models.  A task may mix models simply by issuing prompts that reference a different model id; the framework swaps the backend per call without reinitializing the agent or the chat session.  Beyond bundled providers, \texttt{--endpoint}/\texttt{--header} (and the corresponding \texttt{update\_settings} controls inside the agent) configure any OpenAI-compatible HTTP server---including local models served by, for example, vLLM or llama.cpp---which keeps the bring-your-own-key, local-first posture even for self-hosted backends.  Because keys and prompts travel directly from the developer's machine to the chosen provider, no Sorcar-operated intermediary observes the traffic.


%----------------------------------------------------------------------
\section{Painless Software Engineering with KISS Sorcar}
\label{sec:painless}
%----------------------------------------------------------------------

A central claim of our system is that natural-language interaction can replace manual code inspection and ad hoc scripting for understanding and evolving nontrivial subsystems.  While developing KISS Sorcar, we found a particular workflow to be particularly useful.  In the workflow, we first ask KISS Sorcar to generate a detailed, step-by-step description of a workflow or algorithm we found buggy, and then we ask KISS Sorcar to revise some of the buggy steps in natural language. We illustrate this with a real development session drawn from the project's own history, in which the worktree merge workflow (Section~\ref{sec:worktree-agent}) was first understood and then redesigned entirely through conversational prompts.  The session comprises four consecutive tasks; we reproduce the prompts verbatim and summarize the agent's responses.

\paragraph{Step~1: Understanding the existing workflow.} The developer begins by asking the agent to explain the current post-task git lifecycle:

\begin{prompt}
Can you tell me what happens, step by step, with git in worktree_sorcar_agent.py when a task finishes?
\end{prompt}

\noindent The agent reads the source code and returns a structured summary of the four-phase lifecycle: (1)~during \texttt{run()}, a new branch and worktree are created and the task executes inside the worktree; when the task completes, \emph{nothing is committed or merged}---the result is returned with merge instructions appended and the worktree stays pending; (2)~\texttt{merge()} calls \texttt{\_finalize\_worktree()}, which stages all changes, generates a commit message via the LLM, commits, removes the worktree, checks out the original branch, and runs \texttt{git merge}; (3)~\texttt{discard()} removes the worktree, prunes, checks out the original branch, and deletes the task branch; (4)~in CLI mode an interactive prompt forces the user to choose \texttt{[c]ommit and merge} or \texttt{[d]iscard} before exiting. The agent also notes a key design invariant: nothing auto-merges; auto-commit occurs only at merge/finalize time; all steps are idempotent; and state can be recovered from git on restart.

\paragraph{Step~2: Simplifying the workflow via natural language.} Armed with the workflow description, the developer decides the three-way choice (auto-merge, manual merge, discard) is unnecessarily complex and issues a redesign request:

\begin{prompt}
Can you change worktree_sorcar_agent.py and the extension so that after the agent 
finishes its task, it simply asks "Commit and Merge" or "Discard". When "Commit 
and Merge" is clicked by the user, the agent must commit the changes with a 
generated commit message, merge the branch with the original branch, and delete 
the worktree and the branch associated with the worktree. If the user clicks 
"discard", it must delete the worktree and the branch, and checkout the 
original branch.
\end{prompt}

\noindent The agent modifies six files across Python and TypeScript: it updates \texttt{discard()} to check out the original branch before deleting the task branch, removes the \texttt{manual\_merge()} method entirely, simplifies \texttt{merge\_instructions()} to show only two options, updates the webview UI to replace the three-button toolbar with a two-button ``Commit and Merge or Discard?'' bar, removes the \texttt{manual} action type from the TypeScript type definitions, and removes the corresponding handler from the Python backend.  Three tests for the deleted manual-merge path have been removed, and one routing test has been updated. All 28 worktree tests pass after the change.

\paragraph{Step~3: Investigating unexpected post-merge state.} After testing the simplified flow, the developer notices that files appear in the Source Control panel of VS~Code on the original branch after pressing ``Commit and Merge'' and asks:

\begin{prompt}
After the user presses "Commit and Merge", why do the modified files show up as 
committed in the original branch?
\end{prompt}

\noindent The agent traces the exact execution path through \texttt{squash\_merge\_branch()} and discovers that the implementation deliberately \emph{unstages} the squash-merged changes via \texttt{git reset HEAD} so they appear as uncommitted working-tree modifications in VS~Code's Source Control panel.  The design intention was to let the user review the diff before committing manually.  The agent reports this finding along with the relevant code snippet and a verified end state showing no new commits on the branch.

\paragraph{Step~4: Directing a design change in one sentence.} The developer realizes that code review already happened while the worktree was active and decides the unstage step is unnecessary:

\begin{prompt}
The review is already happening in the worktree branch. You don't need the user 
to review the modified and new files in the original branch. Fix it.
\end{prompt}

\noindent The agent replaces the \texttt{git reset HEAD} call in \texttt{squash\_merge\_branch()} with a conditional \texttt{git commit --no-edit} that uses the auto-generated squash message.  It adds a \texttt{git diff --cached --quiet} guard for the edge case where the merge produces no changes.  One existing test (\texttt{test\_merge\_leaves\_changes\_uncommitted}) is renamed to \texttt{test\_merge\_commits\_changes} and its assertions are updated to verify a clean working tree.  All 104 worktree tests pass.

\medskip
\noindent\textbf{Discussion.} The four-step session illustrates the development loop that our system supports.  The developer never opens a source file, never writes a line of code, and never runs a test manually.  Instead, the entire cycle---understand the workflow, redesign it, investigate an anomaly, direct a fix---happens through natural-language prompts, with the agent handling code reading, multi-file editing, test updates, and verification.  This style of development becomes possible because of the agent hierarchy described in Section~\ref{sec:architecture}: the Worktree Sorcar Agent isolates changes on a branch, the Chat Sorcar Agent preserves conversational context across tasks, and the Relentless Agent automatically continues when the context window is exhausted.

%----------------------------------------------------------------------
\section{Related Work}
\label{sec:related}
%----------------------------------------------------------------------

\textbf{Code-specialized language models.} Beyond the general-purpose LLMs that our system can use as backends, a rich line of work has produced models specifically trained for code. Code Llama~\citep{roziere2023code} fine-tunes Llama~2 for code generation and infilling. StarCoder~\citep{li2023starcoder} trains on permissively licensed code from GitHub with a fill-in-the-middle objective. DeepSeek-Coder~\citep{guo2024deepseek} trains on a 2-trillion-token corpus of code and natural language with a repository-level context window. More recently, frontier models have been optimized specifically for agentic software engineering. Claude Opus~4.6~\citep{anthropic2026opus46} and its successor Opus~4.7~\citep{anthropic2026opus47} advance long-horizon coding and agentic task execution; OpenAI's GPT~5.5~\citep{openai2026gpt55} similarly targets agentic software engineering with strong code generation and tool-use capabilities. Cursor released Composer~2~\citep{cursor2026composer2}, a custom fine-tuned coding model trained with large-scale reinforcement learning. Kimi K2.5~\citep{kimiteam2026k25} is an open-source multimodal agentic model that jointly optimizes text and vision and introduces Agent Swarm for parallel task decomposition. GLM-5.1~\citep{zai2026glm51} is a 754-billion-parameter mixture-of-experts model from Z.ai that sustains autonomous coding over multi-hour sessions, achieving state-of-the-art on SWE-Bench~Pro. Our system is model-agnostic and can leverage any of these models through its pluggable LLM backend, benefiting from advances in code-specialized pre-training without architectural changes.

\textbf{Code generation agents.} SWE-Agent~\citep{yang2024sweagent} and OpenHands~\citep{wang2024openhands} provide LLM-based agents for resolving software engineering tasks such as GitHub issues.  Both use a single-session execution model without automatic continuation. Agentless~\citep{xia2024agentless} takes the opposite approach, showing that a simple two-phase localize-then-repair pipeline without autonomous agent loops can achieve competitive results on SWE-bench~\citep{jimenez2024swebench}.  Devin~\citep{devin2024}, an industrial product marketed as ``the first AI software engineer,'' operates in a sandboxed environment with shell, browser, and editor access.  Claude Code~\citep{claudecode2025} is Anthropic's terminal-based agentic coding tool that gives Claude direct access to a shell, file system, and development tools for multi-step engineering tasks.  OpenAI's Codex~\citep{codex2025} is a cloud-based software engineering agent that executes coding tasks in sandboxed environments with full repository context.  Aider~\citep{aider2023} provides a terminal-based pair programming interface with tight git integration, automatically committing each change.  Our Relentless Agent layer addresses the single-session limitation common to most of these systems, while our worktree isolation provides stronger safety guarantees than per-change commits do.

\textbf{ReAct and tool use.} The ReAct framework~\citep{yao2023react} interleaves reasoning and action. Toolformer~\citep{schick2023toolformer} teaches models to use tools via self-supervised learning.  We use native function calling provided by modern LLM APIs rather than in-context tool descriptions, thereby reducing prompt overhead and improving reliability.

\textbf{Reasoning and planning.} Chain-of-thought prompting~\citep{wei2022chain} demonstrated that eliciting step-by-step reasoning dramatically improves LLM performance on complex tasks.  Tree of Thoughts~\citep{yao2023tree} generalizes this to deliberate search over multiple reasoning paths. Reflexion~\citep{shinn2023reflexion} introduces verbal reinforcement learning, where an agent reflects on failed attempts and produces self-critiques that improve subsequent trials.  Our continuation protocol is conceptually related to Reflexion: failed sub-sessions produce summaries that inform subsequent attempts. However, in our case, we use the summaries to continue the task.

\textbf{Multi-agent software development.} ChatDev~\citep{qian2024chatdev} models the software development process as a conversation between role-playing agents (CEO, CTO, programmer, tester) organized in a waterfall-like pipeline. MetaGPT~\citep{hong2024metagpt} takes a meta-programming approach, encoding standard operating procedures as structured outputs that coordinate specialized agents. AutoGen~\citep{wu2023autogen} provides a general-purpose framework for multi-agent conversation, enabling flexible topologies beyond fixed pipelines.  These systems focus on generating entire applications from scratch.  We take a different approach: rather than simulating an organization of specialists, we use a single agent with broad access to tools and optional parallel sub-agents for embarrassingly parallel sub-tasks, prioritizing practical utility on real-world codebases over role-playing fidelity.

\textbf{Multi-turn autonomous agents.} AutoGPT~\citep{autogpt2023} and BabyAGI~\citep{babyagi2023} implement multi-step autonomous agents.  These systems typically lack budget controls and safe code isolation.  Our layered architecture addresses each of these concerns in a dedicated layer.

\textbf{Software engineering benchmarks.} SWE-bench~\citep{jimenez2024swebench} evaluates agents on real-world GitHub issues drawn from popular Python repositories, requiring the agent to localize and fix bugs given only the issue description. It has become the de facto standard for measuring agent capabilities in software engineering. HumanEval~\citep{chen2021evaluating} and MBPP~\citep{austin2021program} evaluate function-level code generation from docstrings. LiveCodeBench~\citep{jain2024livecodebench} addresses benchmark contamination by continuously collecting fresh competition problems from LeetCode, AtCoder, and CodeForces, and extends evaluation to self-repair, code execution, and test output prediction.  These benchmarks focus on isolated coding problems; We target the broader workflow of multi-file, multi-step software engineering tasks that require tool use, testing, and version control.

\textbf{Agent frameworks and orchestration.} LangChain~\citep{langchain2022} provides modular abstractions for building LLM applications, including agent loops, tool integration, and memory. It focuses on composability and breadth of integrations rather than on the specific concerns of software development.  DSPy~\citep{khattab2024dspy} treats LLM calls as declarative modules whose prompts and few-shot examples are compiled and optimized automatically, enabling systematic improvement of multi-stage pipelines; it targets prompt optimization rather than end-to-end software engineering workflows.  CrewAI~\citep{crewai2024} orchestrates role-playing autonomous agents that collaborate through configurable process models and event-driven flows, emphasizing multi-agent coordination over the layered safety and budget controls that we prioritize.  smolagents~\citep{smolagents2025}, Hugging Face's minimalist agent library, has its CodeAgent write actions as executable Python snippets, reducing step counts by~30\%; however, it does not address repository-level concerns such as git isolation or continuation across sessions.  The OpenAI Agents SDK~\citep{openaiagentssdk2025} offers a provider-agnostic framework for multi-agent workflows with handoffs, guardrails, and tracing, while Google's Agent Development Kit~\citep{googleadk2025} provides an open-source toolkit for building, evaluating, and deploying AI agents with multi-agent orchestration.  Both provide general-purpose agent infrastructure but leave domain-specific concerns to the application layer.  Our architecture is purpose-built for software engineering and general research, with each layer addressing a specific practical concern (budget tracking, continuation, code safety).

\textbf{LLM agent surveys.} Several comprehensive surveys have mapped the rapidly growing landscape of LLM-based agents. \citet{xi2023rise} surveys the design space of LLM agents along three dimensions --- brain (reasoning), perception (input modalities), and action (tool use) --- and catalogs applications across social science, natural science, and engineering. \citet{wang2024survey} propose a systematic framework for autonomous agents built on LLMs, covering profile, memory, planning, and action modules.  Our system can be understood through the lens of these frameworks: the KISS Agent implements the action loop; the Relentless Agent addresses memory and planning concerns through continuation summaries; and the system prompt encodes the profile.

\textbf{IDE integration.} GitHub Copilot~\citep{copilot2021}, Cursor~\citep{cursor2024}, and Windsurf~\citep{windsurf2024} provide AI assistance within editors. Cursor recently released Composer~2~\citep{cursor2026composer2}, a custom fine-tuned coding model trained with large-scale reinforcement learning that achieves 61.7\% on Terminal Bench~2.0. We operate at the level of autonomous multi-step task execution with full tool access and git-level isolation.

\textbf{Recent advances in agentic software engineering.} The field has matured rapidly since late 2025. \citet{hassan2025agentic} articulate the foundational pillars of \emph{Agentic Software Engineering} (SE~3.0), defining a research roadmap that emphasizes trust, controllability, and goal-oriented task decomposition.  \citet{li2025rise} surveys the landscape of autonomous coding agents and characterizes the transition from assistive code completion to full-fledged AI teammates that initiate, plan, and execute development tasks.  We instantiate several of the principles advocated in these roadmaps, including layered controllability (via budget tracking and step limits) and safe isolation (via worktrees).

\citet{wang2025agentic} provides a comprehensive survey of AI agentic programming techniques, cataloging how LLM-based agents decompose goals, interact with compilers and version control systems, and self-correct through feedback loops. \citet{chatlatanagulchai2025readmes} study \emph{agent context files} --- persistent, project-level instruction files that guide agentic coding tools --- and find that high-quality context files significantly improve agent performance.  Our layered system prompt and \texttt{SORCAR.md} override mechanisms are instances of this pattern. \citet{mohsenimofidi2025context} further investigates context engineering for AI agents in open-source software, highlighting how curated context improves agent efficacy on repository-level tasks.

\textbf{Evolving benchmarks for coding agents.} SWE-bench Pro~\citep{deng2025swebenchpro} introduces a substantially more challenging benchmark with 1,865 long-horizon problems drawn from 41 repositories, including commercial codebases, explicitly targeting enterprise-level complexity beyond the original SWE-bench.  These long-horizon tasks --- often requiring multi-file patches and hours of professional developer effort --- directly motivate our continuation mechanism.  \citet{prathifkumar2025swebenchverified} raises the important question of benchmark contamination, showing that overlap between SWE-Bench-Verified problems and LLM training data may inflate scores, suggesting that high benchmark numbers may partly reflect memorization rather than genuine problem-solving ability. \citet{horikawa2025agentic} provides an empirical study of how AI coding agents handle refactoring tasks, finding that while agents can plan and execute complex refactorings, they still struggle with cross-file dependency analysis --- a challenge that our sub-agent parallelism partially addresses.

\textbf{Self-improvement and test-time scaling.} \citet{robeyns2025selfimproving} demonstrates a self-improving coding agent that iteratively refines its own scaffolding code, achieving dramatic benchmark improvements through self-generated optimizations.  We previously experimented with a much lighter-weight variant of this idea via a per-project preferences file that the agent read and updated each session, but removed the mechanism (Section~\ref{sec:self-improvement-prompt}) because the store retained stale project facts that went out of date as the code evolved, leading the agent to reintroduce already-fixed bugs.  \citet{gao2025trae} introduces the Trae Agent, which applies test-time compute scaling to software engineering, dynamically allocating more inference budget to harder problems.  Our budget-tracking mechanism provides a complementary perspective: rather than scaling compute adaptively, we enforce hard budget ceilings while the Relentless Agent ensures maximum progress within those bounds.

\textbf{Community-driven prompt engineering.} Get Shit Done (GSD)~\citep{gsd2025} is a light-weight meta-prompting, context engineering, and spec-driven development system originally created for Claude Code and later extended to other AI coding agents.  GSD addresses \emph{context rot}---the quality degradation that occurs as an LLM fills its context window---by structuring work into phased planning documents and spawning specialized parallel agents for research, planning, execution, and verification.  Its core philosophy rejects enterprise ceremony in favor of directness: ``no enterprise roleplay bullshit,'' as the author puts it.  The sections ``Pre-flight Checks'', ``Deep Work'', and ``Pre-Finish Verification'' of \texttt{SYSTEM.md} were partly inspired by this work.

%----------------------------------------------------------------------
\section{Conclusion}
\label{sec:conclusion}
%----------------------------------------------------------------------

We have shown that a simple layered, single-concern architecture can address the practical challenges of deploying LLM agents for real-world software development.  On Terminal Bench~2.0, a benchmark of 89 diverse terminal-based tasks evaluated across 5 trials each, our system achieves a 62.2\% overall pass rate (277/445), compared with Claude~Code (58\%) and Cursor Composer~2 (61.7\%) on the same benchmark with the same underlying model (Claude Opus~4.6).  These results are achieved without benchmark-specific optimizations, fine-tuning, or reinforcement learning.

We complement the architecture with a system prompt that encodes engineering practices directly into the agent's behavior: read before writing, test before fixing, plan before executing, verify before finishing.  These are not novel insights---they are the practices of careful software engineering, translated into instructions that an LLM can follow.  The evaluation suggests that giving a frontier model the time and tools to validate its own output matters more than model-level customization.

In summary, we showed that a simple agent framework, without sophisticated agent technologies such as trajectory compaction and asynchronous multi-agent orchestration, was sufficient to build KISS Sorcar.  By building KISS Sorcar using itself and matching or exceeding both Cursor and Claude Code, we found that established software engineering techniques and principles are important for building reliable agent systems.

\section*{An Important Advice} We found that researchers are publishing 
lots of papers on AI, and it is hard to keep track of them or validate their claims. We propose to use the following prompt with KISS Sorcar to identify issues in blogs, papers, and code repositories:

\begin{prompt}
Can you read <<url>>, and thoroughly and precisely check for **wrong assumptions**, **cheating**,
**irreproducibility issues**, **fraud**, **potential for cheating in evaluation**
and **security vulnerabilties**?  
Use the internet search extensively and do not believe what people say--verify it yourself.  
Do not hesitate to download the code and run it to validate the results.
For security vulnerabilities, create a POC and test it.  
Generate an HTML report in ./sorcar_reported_frauds/ and open it in the user's default browser.
Thoroughly fact check everything you claim in the report.
\end{prompt}
In the prompt, replace <<url>> with the actual link to a blog, a paper, or a code repository.

\section*{Acknowledgments}

This research is supported in part by gifts from Accenture, Amazon, AMD, Anyscale, Broadcom, Google, IBM, Intel, Intesa Sanpaolo, Lambda, Lightspeed, Mibura, NVIDIA, Samsung SDS, SAP, by the U.S. Department of Energy, Office of Science, Office of Advanced Scientific Computing Research through the X-STACK: Programming Environments for Scientific Computing program (DESC0021982,) and the Defense Advanced Research Projects Agency (DARPA) under Agreement No. HR00112590134.

We would like to thank Marius Momeu for finding critical bugs in the cost computation and in the UI usability, Yogya Mehrotra for testing and fixing bugs in the OpenClaw like third\_party\_agents in KISS Sorcar; Debabrata Dash, Yiwei Hou, Kaiyao Ke, Muxi Lyu, Anoop Mishra, Manish Shetty, Ion Stoica, Shangyin Tan, Hao Wang, and Matei Zaharia for various useful feedback.

Disclaimer: The paper has been edited in part using KISS Sorcar, an AI tool.

%----------------------------------------------------------------------
% References
%----------------------------------------------------------------------


{ \small

\bibliographystyle{plainnat}
\bibliography{kiss_sorcar}
}

\end{document}
